\documentclass[conference]{IEEEtran}
\IEEEoverridecommandlockouts
% The preceding line is only needed to identify funding in the first footnote. If that is unneeded, please comment it out.
\usepackage{cite}
\usepackage{amsmath,amssymb,amsfonts}
% \usepackage{algorithmic}
\usepackage{graphicx}
\usepackage{textcomp}
\usepackage{xcolor}
% \usepackage[margin=1in]{geometry}
\usepackage{amsmath, amssymb, amsthm}
\usepackage{graphicx}
\usepackage{booktabs}
% \usepackage{caption}
\usepackage{subcaption}
\usepackage{algorithm}
% \usepackage{multirow}
\usepackage{algpseudocode}
\usepackage{hyperref}
\usepackage{pgfplots}
\usepackage{pgfplotstable}

\pgfplotsset{compat=1.18}
\def\BibTeX{{\rm B\kern-.05em{\sc i\kern-.025em b}\kern-.08em
    T\kern-.1667em\lower.7ex\hbox{E}\kern-.125emX}}
\begin{document}

\title{TRINITY: A Tri-order Regularized Optimization with Adaptivity}

% \author{\IEEEauthorblockN{1\textsuperscript{st} Given Name Surname}
% \IEEEauthorblockA{\textit{dept. name of organization (of Aff.)} \\
% \textit{name of organization (of Aff.)}\\
% City, Country \\
% email address or ORCID}
% \and
% \IEEEauthorblockN{2\textsuperscript{nd} Given Name Surname}
% \IEEEauthorblockA{\textit{dept. name of organization (of Aff.)} \\
% \textit{name of organization (of Aff.)}\\
% City, Country \\
% email address or ORCID}
% \and
% \IEEEauthorblockN{3\textsuperscript{rd} Given Name Surname}
% \IEEEauthorblockA{\textit{dept. name of organization (of Aff.)} \\
% \textit{name of organization (of Aff.)}\\
% City, Country \\
% email address or ORCID}
% \and
% \IEEEauthorblockN{4\textsuperscript{th} Given Name Surname}
% \IEEEauthorblockA{\textit{dept. name of organization (of Aff.)} \\
% \textit{name of organization (of Aff.)}\\
% City, Country \\
% email address or ORCID}
% \and
% \IEEEauthorblockN{5\textsuperscript{th} Given Name Surname}
% \IEEEauthorblockA{\textit{dept. name of organization (of Aff.)} \\
% \textit{name of organization (of Aff.)}\\
% City, Country \\
% email address or ORCID}
% \and
% \IEEEauthorblockN{6\textsuperscript{th} Given Name Surname}
% \IEEEauthorblockA{\textit{dept. name of organization (of Aff.)} \\
% \textit{name of organization (of Aff.)}\\
% City, Country \\
% email address or ORCID}
% }

\maketitle

\begin{abstract}
The current state of the art in deep learning optimization typically employs zeroth-, first-, and second-order gradient updates. We propose Trinity, a novel optimization framework that unifies these approaches by using zeroth-order curvature sensing for selective natural-gradient preconditioning. Trinity introduces a forward-only second-difference curvature sensor that estimates the effective rank of the Fisher Information Matrix per block at minimal overhead. This sensor output drives a controller that restricts expensive K-FAC inversions to layers where the geometry is highly curved, while an AdamW actuator with magnitude grafting ensures stable scale-safe updates elsewhere. Additionally, the zeroth-order signal acts as a direct saddle-escape actuator. Benchmark evaluations on extreme 100:1 long-tailed label imbalances across diverse architectures demonstrate that Trinity matches the performance of always-on K-FAC at a fraction of the inversion cost.
\end{abstract}

\begin{IEEEkeywords}
Optimizer, K-FAC, Minimization, CNN, Vision Transformer, Gradient Centralization, Preconditioning
\end{IEEEkeywords}



\section{Introduction}
\label{sec:introduction}

First-order optimization methods, though empirically scalable, are fundamentally geometry-blind. Second-order methods, such as natural gradient descent, address this by preconditioning updates with curvature information, but they pay an $\mathcal{O}(n^3)$ cost per block per inversion interval. The standard response to this computational bottleneck is a global schedule, where all layers are inverted uniformly every $\tau$ steps.

However, the value of preconditioning is not uniform across blocks. If a block's Fisher Information Matrix is near-isotropic, the preconditioned gradient $\mathbf{F}^{-1}\nabla L$ is roughly proportional to the raw gradient $\nabla L$, meaning the expensive matrix inversion buys almost nothing. Despite this, such anisotropy is rarely measured during training, and the inversion budget is spent uniformly across the network.

The primary reason this anisotropy remains unmeasured is that traditional diagnostics require Hessian-vector products or full eigendecompositions. These operations are typically as expensive as the preconditioning step they are meant to evaluate, making them impractical for dynamic runtime control.

Our core insight is that two forward passes can provide a second difference rather than a first difference, yielding an unbiased quadratic form $\mathbf{z}^T \mathbf{H} \mathbf{z}$. This operates at the same computational cost as a standard zeroth-order gradient estimate but is significantly more informative, circumventing the $\mathcal{O}(d)$ variance issues that dominate standard backpropagation in differentiable settings.

Building on this, we propose \textbf{Trinity}: a unified framework where zeroth-order (ZO) signals sense the curvature, while second-order (K-FAC) and first-order (AdamW) mechanisms act on it. Furthermore, the ZO component acts directly as a saddle-escape actuator. All three orders are present, each functioning in a structurally defensible role.

Our concrete contributions are:
\begin{enumerate}
    \item \textbf{A forward-only curvature sensor:} Using RDSA second differences, we obtain unbiased estimates of $\mathrm{tr}(\mathbf{H})$ and $\|\mathbf{H}\|_F^2$ per parameter block, yielding a normalized effective rank $\rho \in (0,1]$ at $\sim$6\% overhead.
    \item \textbf{A measurement study of Fisher anisotropy:} We empirically characterize where the metric is actually non-Euclidean across layers and training time under severe label imbalance.
    \item \textbf{Selective natural-gradient preconditioning (Trinity):} A per-block controller spends a bounded K-FAC inversion budget only where the sensor indicates curved geometry, falling back to AdamW elsewhere with scale-safe magnitude grafting.
\end{enumerate}

\section{Related Work}

\subsection{First-Order Adaptive Methods and Structural Regularization}
Stochastic Gradient Descent and its momentum-based variants have historically underpinned deep neural network training. To address sensitivity to hyperparameter tuning and ill-conditioned landscapes, adaptive first-order methods like Adam \cite{kingma2014adam} and RMSProp were introduced. The generalization gap of Adam was fundamentally addressed by Loshchilov and Hutter \cite{loshchilov2019decoupled} with AdamW, which decoupled weight decay from the adaptive gradient updates. Furthering first-order stability, structural constraints applied directly to the gradient tensor, such as Gradient Centralization (GC) \cite{yong2020gradient}, project weight gradients to have a zero mean across non-output dimensions. While these first-order enhancements drastically improve empirical stability, they remain fundamentally oblivious to the underlying curvature of the loss surface.

\subsection{Second-Order Preconditioning and Grafting}
To overcome geometric blindness, second-order optimization explicitly accounts for landscape curvature. Amari's Natural Gradient Descent (NGD) \cite{amari1998natural} adjusts updates using the inverse Fisher Information Matrix (FIM). To achieve tractable scaling, Martens and Grosse \cite{martens2015optimizing} proposed Kronecker-Factored Approximate Curvature (K-FAC) for linear layers, later extended to convolutional layers (KFC) by Grosse and Martens. Extensions of K-FAC have been successfully applied to massive architectures \cite{osawa2019large}. Other recent preconditioning approaches include Shampoo (Gupta et al. 2018), AdaHessian, SOAP (Vyas et al. 2024), and Muon. To safely combine the directional benefits of second-order steps with the robust step sizes of first-order methods, grafting (Agarwal et al.) has emerged as a crucial technique.

\subsection{Zeroth-Order Optimization and Perturbation}
Zeroth-Order (ZO) optimization bypasses analytical backpropagation, constructing trajectories strictly from forward objective evaluations. Methods range from Spall's SPSA \cite{spall1992multivariate} to Nesterov and Spokoiny's Gaussian smoothing estimators \cite{nesterov2017random}. Recently, ZO principles have been integrated into deep learning as landscape-aware regularizers or memory-efficient optimizers, such as MeZO (Malladi et al. 2023). Crucially, Sharpness-Aware Minimization (SAM) (Foret et al. 2021) established the motivation of escaping sharp minima via an extra forward/backward pass. Additionally, perturbed gradient descent (Ge et al. 2015, Jin et al. 2017) demonstrated that injecting random noise allows the optimizer to efficiently escape saddle points.

\subsection{Adaptive Curvature Use}
The classical ancestor of curvature-driven adaptation is the Levenberg-Marquardt (LM) gain ratio, which adapts damping based on the agreement between actual and predicted loss reduction. K-FAC employs an LM-style damping adaptation. We position our approach explicitly against such global, computationally heavy evaluations: our proposed ZO sensing signal is forward-only, dimension-normalized, and evaluated per-block, directly measuring structural anisotropy without requiring expensive matrix inversions or Hessian-vector products.



\section{Background and Notation}
\label{sec:background}

In this section, we establish the mathematical notation, foundational optimization paradigms, and regularized objective functions that form the basis of our work.

\subsection{Problem Setup and Core Notation}
\label{subsec:notation}

We consider a deep neural network parameterized by a weight vector $\boldsymbol{\theta} \in \mathbb{R}^{d}$. The network defines a non-linear mapping from an input token or image $\mathbf{x} \in \mathcal{X}$ to a vector of class logits $\mathbf{z}(\mathbf{x}; \boldsymbol{\theta}) \in \mathbb{R}^{C}$, where $C$ denotes the number of target categories. The conditional probability distribution over the classes is modeled via the softmax operator:
\begin{equation}
p_k(\mathbf{x}; \boldsymbol{\theta}) = \frac{\exp\left(z_k(\mathbf{x}; \boldsymbol{\theta})\right)}{\sum_{j=1}^{C} \exp\left(z_j(\mathbf{x}; \boldsymbol{\theta})\right)}, \qquad \forall k \in \{1, \dots, C\}.
\label{eq:softmax}
\end{equation}

Given a stochastic mini-batch $\mathcal{B} = \{(\mathbf{x}_i, y_i)\}_{i=1}^{B}$ drawn from the training distribution, where $y_i \in \{1, \dots, C\}$ represents the ground-truth label, the empirical risk minimizer seeks to minimize the batch loss function:
\begin{equation}
\mathcal{L}(\boldsymbol{\theta}; \mathcal{B}) = \frac{1}{B}\sum_{i=1}^{B}\ell(\mathbf{x}_i, y_i; \boldsymbol{\theta}),
\label{eq:batch_loss}
\end{equation}
where $\ell(\cdot)$ is the per-sample objective. At optimization step $t$, the standard first-order backpropagated gradient evaluated over the active mini-batch $\mathcal{B}^{(t)}$ is defined as:
\begin{equation}
\mathbf{g}^{(t)} = \nabla_{\boldsymbol{\theta}} \mathcal{L}(\boldsymbol{\theta}^{(t)}; \mathcal{B}^{(t)}).
\label{eq:full_grad}
\end{equation}

For a layer-wise analysis, we define a block partition of the parameters $\boldsymbol{\theta} = \{\boldsymbol{\theta}^{(\ell)}\}$, with corresponding sizes $d_\ell$. We denote the weight matrix and bias vector of layer $\ell$ by $\mathbf{W}^{(\ell)}$ and $\mathbf{b}^{(\ell)}$, respectively. A block is considered eligible for second-order preconditioning if its parameter tensor has dimension $\ge 2$ (e.g., weights of Linear or Convolutional layers) and its Kronecker factors do not exceed a budget $d_{\mathrm{max}}$; we write $B$ for the number of such eligible blocks. One-dimensional blocks, such as normalization gains or standalone biases, together with blocks exceeding $d_{\mathrm{max}}$, are permanently restricted to first-order updates. We denote the matrix of incoming layer activations by $\mathbf{A}^{(\ell)}_{\mathrm{raw}}$ and the matrix of backpropagated pre-activation sensitivities by $\mathbf{S}^{(\ell)}_{\mathrm{raw}}$, each carrying $n$ rows (one per sample, or per sample--spatial-location pair for convolutional layers). The unsubscripted symbols $\mathbf{A}^{(\ell)}, \mathbf{S}^{(\ell)}$ are reserved for the corresponding Kronecker factors defined in Eq.~\eqref{eq:kfac_factors}.

\subsection{First-Order Optimization and Gradient Operations}
\label{subsec:first_order_bg}

\paragraph{Adaptive Moment Estimation (Adam/AdamW)} First-order adaptive optimizers track temporal variations in gradients using exponential moving averages. Standard Adam maintains a first moment vector $\mathbf{m}^{(t)}$ (momentum) and an uncentered second moment vector $\mathbf{v}^{(t)}$ (variance):
\begin{align}
\mathbf{m}^{(t)} &= \beta_1 \mathbf{m}^{(t-1)} + (1-\beta_1)\mathbf{g}^{(t)}, \\
\mathbf{v}^{(t)} &= \beta_2 \mathbf{v}^{(t-1)} + (1-\beta_2)\bigl(\mathbf{g}^{(t)} \odot \mathbf{g}^{(t)}\bigr),
\label{eq:moments}
\end{align}
where $\beta_1, \beta_2 \in (0, 1)$ are decay hyperparameters, and $\odot$ denotes the Hadamard product. Applying bias-correction yields $\hat{\mathbf{m}}^{(t)} = \mathbf{m}^{(t)} / (1-\beta_1^t)$ and $\hat{\mathbf{v}}^{(t)} = \mathbf{v}^{(t)} / (1-\beta_2^t)$. The decoupled weight decay variant, AdamW, separates $L_2$ regularization from these adaptive moment trajectories, yielding the update rule:
\begin{equation}
\boldsymbol{\theta}^{(t+1)} = \boldsymbol{\theta}^{(t)} - \eta^{(t)} \left( \frac{\hat{\mathbf{m}}^{(t)}}{\sqrt{\hat{\mathbf{v}}^{(t)}}+\epsilon} + \lambda_{\mathrm{wd}}\boldsymbol{\theta}^{(t)} \right),
\label{eq:adamw_update}
\end{equation}
where $\eta^{(t)}$ is the learning rate, $\epsilon > 0$ is a stabilization constant, and $\lambda_{\mathrm{wd}}$ is the weight decay coefficient.

\paragraph{Gradient Centralization (GC)} Gradient Centralization acts as a structural constraint directly on the backpropagated gradients of multi-dimensional tensors (e.g., weights of Linear or Conv2d layers). For a weight tensor $\mathbf{W} \in \mathbb{R}^{n \times d_1 \times \cdots \times d_r}$ with computed gradient $\mathbf{G}$, GC projects the gradient matrix by subtracting its mean along the non-output dimensions:
\begin{equation}
\tilde{G}_{i, j_1, \dots, j_r} = G_{i, j_1, \dots, j_r} - \frac{1}{\prod_{k=1}^r d_k} \sum_{k_1,\dots,k_r} G_{i, k_1, \dots, k_r}.
\label{eq:gradient_centralization}
\end{equation}

\subsection{Curvature-Aware Optimization via K-FAC}
\label{subsec:kfac_bg}

Natural Gradient Descent (NGD) scales the gradient by the inverse of the Fisher Information Matrix (FIM), denoted as $\mathbf{F}(\boldsymbol{\theta})$:
\begin{equation}
\Delta \boldsymbol{\theta}_{\mathrm{NGD}} = -\eta \mathbf{F}(\boldsymbol{\theta})^{-1} \nabla_{\boldsymbol{\theta}} \mathcal{L}(\boldsymbol{\theta}).
\label{eq:ngd}
\end{equation}
As explicitly calculating and inverting the full $d \times d$ matrix $\mathbf{F}$ is computationally expensive, Kronecker-Factored Approximate Curvature (K-FAC) introduces a block-diagonal approximation. For a structural layer $\ell$, the corresponding diagonal block $\mathbf{F}^{(\ell)}$ is factored into the Kronecker product ($\otimes$) of two significantly smaller covariance matrices:
\begin{equation}
\begin{aligned}
\mathbf{F}^{(\ell)} &\approx \mathbf{A}^{(\ell)} \otimes \mathbf{S}^{(\ell)}, \\
\text{where } \mathbf{A}^{(\ell)} &= \mathbb{E}\bigl[\mathbf{A}^{(\ell)\top}_{\mathrm{raw}}\mathbf{A}^{(\ell)}_{\mathrm{raw}}\bigr], \\
\mathbf{S}^{(\ell)} &= \mathbb{E}\bigl[\mathbf{S}^{(\ell)\top}_{\mathrm{raw}}\mathbf{S}^{(\ell)}_{\mathrm{raw}}\bigr].
\end{aligned}
\label{eq:kfac_factors}
\end{equation}

On a mini-batch the expectations in Eq.~\eqref{eq:kfac_factors} are replaced by their batch-averaged counterparts $\tfrac{1}{n}\mathbf{A}^{(\ell)\top}_{\mathrm{raw}}\mathbf{A}^{(\ell)}_{\mathrm{raw}}$ and $\tfrac{1}{n}\mathbf{S}^{(\ell)\top}_{\mathrm{raw}}\mathbf{S}^{(\ell)}_{\mathrm{raw}}$, where the $1/n$ normalization keeps the factors and hence the damping scale applied to them invariant to batch size and spatial resolution. These empirical factors are tracked by an exponential moving average with decay $\rho_{\mathrm{kf}} \in (0,1)$ (distinct from the effective-rank ratio $\hat{\rho}$ of Section~\ref{subsec:zo_sensor}):
\begin{equation}
\label{eq:kfac_ema}
\begin{aligned}
\hat{\mathbf{A}}^{(\ell)}_t &= \rho_{\mathrm{kf}} \hat{\mathbf{A}}^{(\ell)}_{t-1} + (1-\rho_{\mathrm{kf}})\mathbf{A}^{(\ell)}_{\text{emp},t}, \\
\hat{\mathbf{S}}^{(\ell)}_t &= \rho_{\mathrm{kf}} \hat{\mathbf{S}}^{(\ell)}_{t-1} + (1-\rho_{\mathrm{kf}})\mathbf{S}^{(\ell)}_{\text{emp},t},
\end{aligned}
\end{equation}
initialized at each block's first observation rather than from an identity prior.

For numerical stability we adopt the standard trace-balanced Tikhonov damping of~\cite{martens2015optimizing}, which splits a damping coefficient across the two factors via a trace-ratio scalar $\pi^{(\ell)}$ before inversion. The resulting damped inverses $\mathbf{A}^{(\ell)}_{\mathrm{inv}}, \mathbf{S}^{(\ell)}_{\mathrm{inv}}$ are applied to the combined weight--bias gradient $\mathbf{C}^{(\ell)} = [\mathbf{G}^{(\ell)}\;\mathbf{g}_b^{(\ell)}]$ as $\mathbf{S}^{(\ell)}_{\mathrm{inv}}\mathbf{C}^{(\ell)}\mathbf{A}^{(\ell)}_{\mathrm{inv}}$, which is the Kronecker-factored realization of Eq.~\eqref{eq:ngd}; the exact damping and preconditioning expressions appear in Algorithm~\ref{alg:trinity}.

\subsection{Second-Difference Curvature Sensing}
\label{subsec:rdsa_bg}
Rather than estimating the gradient, symmetric directional evaluations can be used to construct a second-difference identity that measures local curvature. For a parameter block $\boldsymbol{\theta}^{(\ell)}$, let $\mathbf{z} \sim \mathcal{N}(\mathbf{0}, \mathbf{I}_{d_\ell})$ be a random Gaussian perturbation supported on block $\ell$ alone, all other blocks being held fixed. By evaluating the loss at $\boldsymbol{\theta}^{(\ell)} \pm \epsilon_{\mathrm{Z}} \mathbf{z}$, the second difference is defined as:
\begin{equation}
c = \frac{\mathcal{L}(\boldsymbol{\theta}^{(\ell)} + \epsilon_{\mathrm{Z}} \mathbf{z}; \mathcal{B}) - 2\mathcal{L}(\boldsymbol{\theta}^{(\ell)}; \mathcal{B}) + \mathcal{L}(\boldsymbol{\theta}^{(\ell)} - \epsilon_{\mathrm{Z}} \mathbf{z}; \mathcal{B})}{\epsilon_{\mathrm{Z}}^2}.
\label{eq:second_diff}
\end{equation}
Expanding both perturbed evaluations about $\boldsymbol{\theta}^{(\ell)}$, the odd-order terms cancel and the even-order terms survive, giving
\begin{equation}
c = \mathbf{z}^\top \mathbf{H}^{(\ell)} \mathbf{z} \;+\; \frac{\epsilon_{\mathrm{Z}}^2}{12}\,D^4\mathcal{L}[\mathbf{z}]^4 \;+\; \mathcal{O}(\epsilon_{\mathrm{Z}}^4),
\label{eq:second_diff_expansion}
\end{equation}
where $\mathbf{H}^{(\ell)}$ is the local Hessian of the loss with respect to block $\ell$ and $D^4\mathcal{L}[\mathbf{z}]^4$ denotes the fourth directional derivative along $\mathbf{z}$. The second difference is therefore a \emph{second-order accurate} estimate of the quadratic form $\mathbf{z}^\top \mathbf{H}^{(\ell)} \mathbf{z}$: exact for quadratic objectives, and carrying an $\mathcal{O}(\epsilon_{\mathrm{Z}}^2)$ truncation bias otherwise. Unbiasedness enters one level higher, through the randomization over $\mathbf{z}$ rather than through the finite difference itself; this is what the moment estimators of Section~\ref{subsec:zo_sensor} exploit, and it is why $\epsilon_{\mathrm{Z}}$ cannot simply be driven to zero.

\section{Proposed Approach: Trinity}
\label{sec:method}

\subsection{Design Principle}
\label{subsec:design_principle}
The core principle of Trinity is a strict reassignment of optimization roles: zeroth-order signals sense the geometry, while first- and second-order methods act on it. A purely additive tri-order composition suffers from $\mathcal{O}(d)$ variance in the ZO gradient estimator and biases the descent direction through layer-wise rescaling. Instead, Trinity structurally decouples the sensing mechanism from the update engine.

\subsection{ZO Curvature Sensor}
\label{subsec:zo_sensor}
For a block $\ell$, Trinity draws $m$ probes $\mathbf{z}_1, \dots, \mathbf{z}_m \sim \mathcal{N}(\mathbf{0}, \mathbf{I}_{d_\ell})$ and evaluates Eq.~\eqref{eq:second_diff} for each, giving $c_i \approx \mathbf{z}_i^\top \mathbf{H}^{(\ell)} \mathbf{z}_i$. The sensor rests on two exact moment identities for Gaussian quadratic forms, $\mathbb{E}[\mathbf{z}^\top\mathbf{H}\mathbf{z}] = \mathrm{tr}(\mathbf{H})$ and $\mathrm{Var}[\mathbf{z}^\top\mathbf{H}\mathbf{z}] = 2\,\mathrm{tr}(\mathbf{H}^2) = 2\|\mathbf{H}\|_F^2$: writing $\bar{c}$ and $s^2$ for the sample mean and unbiased sample variance of $\{c_i\}$, two forward-only moments recover both spectral summaries, $\mathrm{tr}(\mathbf{H}^{(\ell)}) \approx \bar{c}$ and $\|\mathbf{H}^{(\ell)}\|_F^2 \approx s^2/2$. These determine the \emph{effective rank} $r^{(\ell)}_{\mathrm{eff}} = (\mathrm{tr}\,\mathbf{H}^{(\ell)})^2 / \|\mathbf{H}^{(\ell)}\|_F^2 \in [1, d_\ell]$, which attains $d_\ell$ for an isotropic block and $1$ for a rank-one block. Since $\mathbb{E}[\bar{c}^2] = (\mathrm{tr}\,\mathbf{H})^2 + \mathrm{Var}(c)/m$, the plug-in $\bar{c}^2$ over-estimates $(\mathrm{tr}\,\mathbf{H})^2$; subtracting $s^2/m$ removes exactly this term, yielding the \emph{debiased} estimator, its dimension-normalized form, and the companion negative-curvature fraction:
\begin{equation}
\begin{aligned}
\hat{r}_{\mathrm{eff}} = \frac{2\max(\bar{c}^2 - s^2/m,\, 0)}{s^2 + \delta},
\\
\quad
\hat{\rho} = \mathrm{clip}\!\left(\frac{\hat{r}_{\mathrm{eff}}}{d_\ell},\, \frac{1}{d_\ell},\, 1\right),
\quad
\hat{\nu} = \frac{|\{i : c_i < 0\}|}{m},
\label{eq:rho_hat}
\end{aligned}
\end{equation}
so that $\hat{\rho} \in [1/d_\ell, 1]$, with $\delta > 0$ guarding the degenerate case $s^2 \to 0$. Both statistics cost $2m+1$ forward passes and no backward pass; $\hat{\nu}$ is the signal consumed by the escape actuator of Section~\ref{subsec:escape}.

Three correctness requirements govern the validity of this sensor:
\begin{enumerate}
    \item \textbf{Shared Minibatch:} The same minibatch must serve all $m$ probes and the reference evaluation. Resampling injects data-sampling variance into $s^2$, which Eq.~\eqref{eq:rho_hat} would read as curvature anisotropy.
    \item \textbf{High-Precision Accumulation:} Probe losses must be accumulated in fp32 or fp64. The numerator of Eq.~\eqref{eq:second_diff} is a cancelling sum of surviving magnitude $\mathcal{O}(\epsilon_{\mathrm{Z}}^2)$, and the $1/\epsilon_{\mathrm{Z}}^2$ division amplifies the residual roundoff accordingly; in bf16/fp16 the result is roundoff-dominated.
    \item \textbf{Perturbation Radius:} $\epsilon_{\mathrm{Z}}$ must remain macroscopic, as the two error sources oppose one another: truncation bias grows as $\mathcal{O}(\epsilon_{\mathrm{Z}}^2)$ (Eq.~\eqref{eq:second_diff_expansion}), while cancellation error contributes $\mathcal{O}(u|\mathcal{L}|/\epsilon_{\mathrm{Z}}^2)$ for unit roundoff $u$. Balancing the two gives $\epsilon_{\mathrm{Z}}^\star \sim u^{1/4}$, i.e.\ $\approx 10^{-2}$ in fp32, the value used throughout; shrinking $\epsilon_{\mathrm{Z}}$ in pursuit of a tighter derivative approximation is actively harmful and yields noise difficult to distinguish from a functioning sensor.
\end{enumerate}
To protect the Fisher statistics of Eq.~\eqref{eq:kfac_ema} from the sensor's off-path forward passes, a collection flag ($\mathcal{F}_{\mathrm{collect}}$) gates the activation-capturing hooks during sensing, so that $\hat{\mathbf{A}}^{(\ell)}, \hat{\mathbf{S}}^{(\ell)}$ reflect only on-trajectory statistics.

\subsection{Information-Geometric Interpretation}
\label{subsec:info_geom}
The normalized effective rank characterizes isotropy exactly, rather than merely correlating with it. For symmetric $\mathbf{H}$ with eigenvalues $\{\lambda_i\}_{i=1}^{d}$, Cauchy--Schwarz gives $(\mathrm{tr}\,\mathbf{H})^2 = (\sum_i \lambda_i)^2 \le d \sum_i \lambda_i^2 = d\|\mathbf{H}\|_F^2$, with equality \emph{iff} all eigenvalues coincide; hence $\hat{\rho} \to 1$ certifies $\mathbf{H}^{(\ell)} \to c\mathbf{I}$ and $\hat{\rho} \ll 1$ certifies a strongly anisotropic spectrum. Transferring this from the Hessian to the Fisher invokes the standard Gauss--Newton correspondence: for the softmax cross-entropy objective of Eqs.~\eqref{eq:softmax}--\eqref{eq:batch_loss}, the Fisher coincides with the generalized Gauss--Newton matrix, which agrees with the Hessian up to a residual-weighted term that vanishes as the model fits. An isotropic block Hessian then implies $\mathbf{F}^{(\ell)} \approx c\mathbf{I}$, in which regime the natural gradient degenerates to the Euclidean gradient ($\mathbf{F}^{-1}\nabla \mathcal{L} \propto \nabla \mathcal{L}$) and the $\mathcal{O}(n^3)$ inversion purchases nothing, whereas $\hat{\rho} \ll 1$ signals geometry where the Kronecker structure of K-FAC carries genuine value. This motivates pairing the sensor with K-FAC specifically: it measures the same spectral anisotropy that the actuator exploits.

Two caveats delimit the sensor's scope. The range $r^{(\ell)}_{\mathrm{eff}} \in [1, d_\ell]$ presumes a positive semi-definite spectrum in an indefinite region, cancellation within $\mathrm{tr}\,\mathbf{H}^{(\ell)}$ drives $\hat{r}_{\mathrm{eff}}$ toward zero, so $\hat{\rho}$ saturates at its floor $1/d_\ell$ and there reports \emph{indefiniteness} rather than anisotropy. Moreover, K-FAC's Fisher approximation is positive semi-definite by construction and is therefore blind to the negative spectrum that $\hat{\nu}$ detects. Trinity separates these roles explicitly rather than conflating them: $\hat{\rho}$ routes a block to the second-order actuator, $\hat{\nu}$ to the escape actuator of Section~\ref{subsec:escape}.

\subsection{Round-Robin Scheduling}
\label{subsec:round_robin}
Sensing all $B$ blocks per event would require $(2m+1)B$ forward passes, which is prohibitive moreover, per-block probes cannot share forward passes, since a joint perturbation measures cross-block Hessian terms rather than the diagonal block of interest. Trinity therefore visits only $n_p$ blocks per sensing interval $\tau_{\mathrm{Z}}$, cycling round-robin through the eligible set. Charging a training step at three forward-equivalents (one forward, two for the backward), the amortized sensing overhead
\begin{equation}
\mathcal{O}_{\mathrm{sense}} = \frac{(2m+1)\,n_p}{3\,\tau_{\mathrm{Z}}}
\label{eq:sense_overhead}
\end{equation}
depends only on the sensing triple $(m, n_p, \tau_{\mathrm{Z}})$, independent of $B$ and of the parameter count: at $m=4$, $n_p=2$, $\tau_{\mathrm{Z}}=100$ it evaluates to $6\%$. The cost is latency rather than throughput: a full sweep takes $\lceil B/n_p \rceil \tau_{\mathrm{Z}}$ steps, and since $\tilde{\rho}^{(\ell)}$ is initialized to $1$ and moves only when its block is probed, no block is promoted before it is first visited. Sensing latency therefore grows as $\mathcal{O}(B/n_p)$, so $n_p$ must scale with the block count to keep the controller responsive on wider architectures. We report $(m, n_p, \tau_{\mathrm{Z}})$ and the resulting $\mathcal{O}_{\mathrm{sense}}$ per architecture in Section~\ref{sec:experimental_setup}.

\subsection{Controller}
\label{subsec:controller}
The decision to activate K-FAC preconditioning for a given block is governed by a stateful per-block controller, invoked once per sensing event. It maintains a smoothed estimate, updated whenever the block is probed,
\begin{equation}
\tilde{\rho}^{(\ell)} \leftarrow \alpha_{\mathrm{ema}}\,\tilde{\rho}^{(\ell)} + (1 - \alpha_{\mathrm{ema}})\,\hat{\rho}^{(\ell)},
\qquad \tilde{\rho}^{(\ell)} \text{ initialized to } 1,
\label{eq:rho_ema}
\end{equation}
and admits blocks to a candidate set under asymmetric hysteresis: a block currently in FO mode is admitted only if $\tilde{\rho}^{(\ell)} < \rho_{\mathrm{lo}}$, whereas a block already in SO mode is retained while $\tilde{\rho}^{(\ell)} \le \rho_{\mathrm{hi}}$. With $\rho_{\mathrm{lo}} \le \rho_{\mathrm{hi}}$, the band $[\rho_{\mathrm{lo}}, \rho_{\mathrm{hi}}]$ forms a dead zone in which no block changes mode. 

Finally, a global budget caps concurrent inversions: should the candidate set exceed $\lfloor \beta_{\mathrm{max}} B \rfloor$, only the $\lfloor \beta_{\mathrm{max}} B \rfloor$ candidates with the smallest $\tilde{\rho}$ are retained, where $B$ counts SO-eligible blocks only. Because $\tilde{\rho}^{(\ell)}$ is initialized to $1$ and is revised only when block $\ell$ is probed, an unvisited block can never be promoted; the round-robin coverage rate of Section~\ref{subsec:round_robin} therefore rate-limits the onset of preconditioning, and the budget binds only once enough blocks have been sensed to compete for it.

\subsection{Grafting}
\label{subsec:grafting}
First-order (FO) and second-order (SO) modes inherently operate on different step-size scales. A naive switch between modes would induce a sudden jump in the effective step size by an unknown, layer-dependent factor. To ensure a scale-safe transition, Trinity employs magnitude grafting. A shadow AdamW state is permanently maintained for every block to track FO moments. When a block is updated, the \textit{direction} is taken from the active mode (SO if preconditioned, FO otherwise), but the \textit{magnitude} (Frobenius norm) is strictly grafted from the shadow AdamW state. Writing $\boldsymbol{\Delta}_{\mathrm{ad}}$ for the shadow AdamW step and $\boldsymbol{\Delta}_{\mathrm{act}}$ for the active-mode step, the applied update is
\begin{equation}
\boldsymbol{\Delta} = \|\boldsymbol{\Delta}_{\mathrm{ad}}\|_F \cdot \frac{\boldsymbol{\Delta}_{\mathrm{act}}}{\|\boldsymbol{\Delta}_{\mathrm{act}}\|_F + \delta},
\label{eq:grafting}
\end{equation}
so that $\|\boldsymbol{\Delta}\|_F \approx \|\boldsymbol{\Delta}_{\mathrm{ad}}\|_F$ while $\boldsymbol{\Delta}$ remains parallel to $\boldsymbol{\Delta}_{\mathrm{act}}$. Both norms are taken per block, jointly over the concatenated weight and bias $[\mathbf{W}^{(\ell)}\;\mathbf{b}^{(\ell)}]$, so the graft is invariant to how a block's parameters are partitioned across tensors. In FO mode $\boldsymbol{\Delta}_{\mathrm{act}} = \boldsymbol{\Delta}_{\mathrm{ad}}$ and Eq.~\eqref{eq:grafting} reduces to the identity; grafting is thus active only on preconditioned blocks, and the FO path is left numerically untouched.

\subsection{ZO as Escape Actuator}
\label{subsec:escape}
In addition to guiding the controller, the ZO sensor functions directly as a saddle-escape actuator, reusing probes that have already been paid for. A high negative-curvature fraction paired with a vanishing gradient is the signature of a saddle specifically: at a local minimum the block Hessian is positive semi-definite and $\hat{\nu} \approx 0$, so the conjunction of the two conditions excludes minima by construction. A block is perturbed when
\begin{equation}
\tilde{\nu}^{(\ell)} \ge \nu_{\mathrm{thr}}
\quad\text{and}\quad
\|\mathbf{g}^{(\ell)}\|_2 \le \gamma_g\,\bar{g}^{(\ell)}
\quad\text{and}\quad
\mathrm{lock}^{(\ell)} = 0,
\label{eq:escape_trigger}
\end{equation}
where $\bar{g}^{(\ell)}$ is an exponential moving average of the block's own gradient norm, so that the second test asks whether the gradient has collapsed relative to its recent scale rather than against an absolute threshold a scale-free criterion that transfers across layers of very different gradient magnitudes. The perturbation is a single isotropic step scaled by the block's parameter norm,
\begin{equation}
\boldsymbol{\theta}^{(\ell)} \leftarrow \boldsymbol{\theta}^{(\ell)} + r_{\mathrm{esc}}\|\boldsymbol{\theta}^{(\ell)}\|_2\,\mathbf{u},
\qquad \mathbf{u} \sim \mathrm{Uniform}\bigl(\mathcal{S}^{d_\ell - 1}\bigr),
\label{eq:escape_step}
\end{equation}
after which the block's momentum buffers are halved they describe a trajectory the jump has invalidated and a lockout counter is set to $\tau_{\mathrm{esc}}$ steps to prevent immediate re-triggering.


% \begin{algorithm*}[!t]
% \footnotesize
% \caption{Trinity Optimization Framework}
% \label{alg:trinity}
% \begin{algorithmic}[1]
% \Require Parameters $\boldsymbol{\theta}^{(0)}$; learning rate schedule $\eta^{(t)}$; loss function $L(\cdot)$; blocks $\{\ell\}$, sizes $d_\ell$
% \Require AdamW: $\beta_1, \beta_2, \epsilon_a, \lambda_{\mathrm{wd}}$
% \Require Sensor: $\tau_{\mathrm{Z}}, m, \epsilon_{\mathrm{Z}}, n_p$
% \Require Control: $\rho_{\mathrm{lo}}, \rho_{\mathrm{hi}}, \alpha_{\mathrm{ema}}, \tau_{\mathrm{dwell}}, \beta_{\mathrm{max}}$
% \Require K-FAC: $\tau_{\mathrm{stat}}, \tau_{\mathrm{inv}}, \rho_{\mathrm{kf}}, \lambda_{\mathrm{damp}}, d_{\mathrm{max}}$
% \Require Escape: $\nu_{\mathrm{thr}}, \gamma_g, r_{\mathrm{esc}}, \tau_{\mathrm{esc}}$
% \State \textbf{Initialize:} For all $\ell$: $\mathbf{M}, \mathbf{V} \leftarrow 0$ (SO branch); $\mathbf{M}_{\mathrm{ad}}, \mathbf{V}_{\mathrm{ad}} \leftarrow 0$ (shadow AdamW); $\hat{\mathbf{A}}, \hat{\mathbf{S}}, \mathbf{A}_{\mathrm{inv}}, \mathbf{S}_{\mathrm{inv}} \leftarrow \mathbf{I}$
% \State \textbf{Initialize:} For all $\ell$: $\mathrm{mode} \leftarrow \mathrm{FO}, \tilde{\rho} \leftarrow 1, \tilde{\nu} \leftarrow 0, \mathrm{dwell} \leftarrow 0, \mathrm{lock} \leftarrow 0, \bar{g} \leftarrow 0$; $\mathcal{F}_{\mathrm{collect}} \leftarrow \mathrm{True}; t \leftarrow 0$
% \While{not converged}
%     \State $t \leftarrow t + 1$; sample minibatch $\mathcal{B}^{(t)}$
%     \State $L_{\mathrm{base}} \leftarrow L(\boldsymbol{\theta}^{(t)}; \mathcal{B}^{(t)})$; $\mathbf{g} \leftarrow \nabla_{\boldsymbol{\theta}} L_{\mathrm{base}}$ \Comment{Hooks cache $\mathbf{A}_{\mathrm{raw}}, \mathbf{S}_{\mathrm{raw}}$}
%     \If{use\_GC}
%         \State $\mathbf{g}^{(\ell)} \leftarrow \mathrm{CENTRALIZE}(\mathbf{g}^{(\ell)})$ for all $\ell$ with $\mathrm{ndim} \ge 2$
%     \EndIf

%     \Statex \textbf{Phase 1: ZO Sensing (every $\tau_{\mathrm{Z}}$ steps)}
%     \If{$t \bmod \tau_{\mathrm{Z}} == 0$}
%         \State $\mathcal{F}_{\mathrm{collect}} \leftarrow \mathrm{False}$ \Comment{Freeze Fisher statistics}
%         \State $S_t \leftarrow$ next $n_p$ eligible blocks (round-robin)
%         \For{$\ell \in S_t$}
%             \State $(\hat{\rho}, \hat{\nu}) \leftarrow \mathrm{ZO\_PROBE}(\boldsymbol{\theta}^{(t)}, \mathcal{B}^{(t)}, \ell, m, \epsilon_{\mathrm{Z}})$ \Comment{Second-difference sensing}
%             \State $\tilde{\rho}^{(\ell)} \leftarrow \alpha_{\mathrm{ema}} \tilde{\rho}^{(\ell)} + (1-\alpha_{\mathrm{ema}}) \hat{\rho}$; $\tilde{\nu}^{(\ell)} \leftarrow \hat{\nu}$
%         \EndFor
%         \State $\mathcal{F}_{\mathrm{collect}} \leftarrow \mathrm{True}$
%         \State $\{\mathrm{mode}^{(\ell)}\} \leftarrow \mathrm{CONTROLLER}(\{\tilde{\rho}\}, \{\mathrm{dwell}\}, \beta_{\mathrm{max}})$
%     \EndIf

%     \Statex \textbf{Phase 2: ZO Actuator (Saddle Escape)}
%     \For{$\ell \in S_t$ sensed at this step}
%         \If{$\tilde{\nu}^{(\ell)} \ge \nu_{\mathrm{thr}}$ and $\|\mathbf{g}^{(\ell)}\|_2 \le \gamma_g \bar{g}^{(\ell)}$ and $\mathrm{lock}^{(\ell)} == 0$}
%             \State $\mathbf{u} \sim \mathrm{Uniform}(\mathcal{S}^{d_\ell - 1})$
%             \State $\boldsymbol{\theta}^{(\ell)} \leftarrow \boldsymbol{\theta}^{(\ell)} + r_{\mathrm{esc}} \|\boldsymbol{\theta}^{(\ell)}\|_2 \mathbf{u}$ \Comment{Perturbed GD}
%             \State $\mathbf{M}^{(\ell)} \leftarrow \frac{1}{2}\mathbf{M}^{(\ell)}$; $\mathbf{M}_{\mathrm{ad}}^{(\ell)} \leftarrow \frac{1}{2}\mathbf{M}_{\mathrm{ad}}^{(\ell)}$ \Comment{Damp stale momentum}
%             \State $\mathrm{lock}^{(\ell)} \leftarrow \tau_{\mathrm{esc}}$
%         \EndIf
%     \EndFor
%     \For{all $\ell$} $\mathrm{lock}^{(\ell)} \leftarrow \max(\mathrm{lock}^{(\ell)} - 1, 0)$; $\bar{g}^{(\ell)} \leftarrow \mathrm{EMA}(\|\mathbf{g}^{(\ell)}\|_2)$ \EndFor

%     \Statex \textbf{Phase 3: Fisher Factors and Inverses}
%     \If{$t \bmod \tau_{\mathrm{stat}} == 0$}
%         \For{all eligible $\ell$}
%             \State $\hat{\mathbf{A}}^{(\ell)} \leftarrow \rho_{\mathrm{kf}}\hat{\mathbf{A}}^{(\ell)} + (1-\rho_{\mathrm{kf}})\mathbf{A}_{\mathrm{raw}}^T \mathbf{A}_{\mathrm{raw}}$; $\hat{\mathbf{S}}^{(\ell)} \leftarrow \rho_{\mathrm{kf}}\hat{\mathbf{S}}^{(\ell)} + (1-\rho_{\mathrm{kf}})\mathbf{S}_{\mathrm{raw}}^T \mathbf{S}_{\mathrm{raw}}$
%         \EndFor
%     \EndIf
%     \If{$t \bmod \tau_{\mathrm{inv}} == 0$}
%         \For{$\ell$ with $\mathrm{mode}^{(\ell)} == \mathrm{SO}$ and $\dim(\ell) \le d_{\mathrm{max}}$}
%             \State $\pi \leftarrow \sqrt{(\mathrm{tr}(\hat{\mathbf{A}}) \cdot \dim(\hat{\mathbf{S}})) / (\mathrm{tr}(\hat{\mathbf{S}}) \cdot \dim(\hat{\mathbf{A}}) + \delta)}$
%             \State $\mathbf{A}_{\mathrm{inv}}^{(\ell)} \leftarrow (\hat{\mathbf{A}} + \sqrt{\lambda_{\mathrm{damp}}} \pi \mathbf{I})^{-1}$; $\mathbf{S}_{\mathrm{inv}}^{(\ell)} \leftarrow (\hat{\mathbf{S}} + (\sqrt{\lambda_{\mathrm{damp}}}/\pi) \mathbf{I})^{-1}$
%         \EndFor
%     \EndIf

%     \Statex \textbf{Phase 4: Update Engine with Grafting}
%     \For{all $\ell$}
%         \State $\mathbf{M}_{\mathrm{ad}} \leftarrow \beta_1 \mathbf{M}_{\mathrm{ad}} + (1-\beta_1)\mathbf{g}^{(\ell)}$; $\mathbf{V}_{\mathrm{ad}} \leftarrow \beta_2 \mathbf{V}_{\mathrm{ad}} + (1-\beta_2)(\mathbf{g}^{(\ell)} \odot \mathbf{g}^{(\ell)})$
%         \State $\boldsymbol{\Delta}_{\mathrm{ad}} \leftarrow (\mathbf{M}_{\mathrm{ad}} / (1-\beta_1^t)) \oslash (\sqrt{\mathbf{V}_{\mathrm{ad}} / (1-\beta_2^t)} + \epsilon_a)$ \Comment{Shadow AdamW magnitude}
%         \If{$\mathrm{mode}^{(\ell)} == \mathrm{SO}$}
%             \State $\mathbf{C} \leftarrow [\mathbf{G}^{(\ell)} \;\; \mathbf{g}_b^{(\ell)}]$
%             \State $\mathbf{G}_{\mathrm{pre}}, \mathbf{g}_{b,\mathrm{pre}} \leftarrow \mathrm{Split}(\mathbf{S}_{\mathrm{inv}}^{(\ell)} \mathbf{C} \mathbf{A}_{\mathrm{inv}}^{(\ell)})$
%             \State $\mathbf{M} \leftarrow \beta_1 \mathbf{M} + (1-\beta_1)\mathbf{G}_{\mathrm{pre}}$; $\mathbf{V} \leftarrow \beta_2 \mathbf{V} + (1-\beta_2)(\mathbf{G}_{\mathrm{pre}} \odot \mathbf{G}_{\mathrm{pre}})$
%             \State $\boldsymbol{\Delta}_{\mathrm{act}} \leftarrow (\mathbf{M} / (1-\beta_1^t)) \oslash (\sqrt{\mathbf{V} / (1-\beta_2^t)} + \epsilon_a)$
%         \Else
%             \State $\boldsymbol{\Delta}_{\mathrm{act}} \leftarrow \boldsymbol{\Delta}_{\mathrm{ad}}$; $\mathbf{M} \leftarrow \mathbf{M}_{\mathrm{ad}}$; $\mathbf{V} \leftarrow \mathbf{V}_{\mathrm{ad}}$ \Comment{Keep SO state warm}
%         \EndIf
%         \State $\boldsymbol{\Delta} \leftarrow \|\boldsymbol{\Delta}_{\mathrm{ad}}\|_F \cdot \boldsymbol{\Delta}_{\mathrm{act}} / (\|\boldsymbol{\Delta}_{\mathrm{act}}\|_F + \delta)$ \Comment{GRAFTING}
%         \State $\boldsymbol{\theta}^{(\ell)} \leftarrow \boldsymbol{\theta}^{(\ell)} - \eta^{(t)}(\boldsymbol{\Delta} + \lambda_{\mathrm{wd}}\boldsymbol{\theta}^{(\ell)})$
%     \EndFor
% \EndWhile
% \end{algorithmic}
% \end{algorithm*}


\begin{algorithm*}[!t]
\footnotesize
\caption{Trinity Optimization Framework}
\label{alg:trinity}
\begin{algorithmic}[1]
\Require $\boldsymbol{\theta}^{(0)}$; LR schedule $\eta^{(t)}$; loss $L(\cdot)$; $N$ blocks $\{\ell\}$ of size $d_\ell$, of which $B$ are SO-eligible ($d_\ell \le d_{\mathrm{max}}$)
\Require AdamW $\beta_1, \beta_2, \epsilon_a, \lambda_{\mathrm{wd}}$; Sensor $\tau_{\mathrm{Z}}, m, \epsilon_{\mathrm{Z}}, n_p$; Control $\rho_{\mathrm{lo}}, \rho_{\mathrm{hi}}, \alpha_{\mathrm{ema}}, \tau_{\mathrm{dwell}}, \beta_{\mathrm{max}}$; K-FAC $\tau_{\mathrm{stat}}, \tau_{\mathrm{inv}}, \rho_{\mathrm{kf}}, \lambda_{\mathrm{damp}}, d_{\mathrm{max}}$; Escape $\nu_{\mathrm{thr}}, \gamma_g, r_{\mathrm{esc}}, \tau_{\mathrm{esc}}$
\State \textbf{Init:} $\forall \ell$: $\mathbf{M}, \mathbf{V}, \mathbf{M}_{\mathrm{ad}}, \mathbf{V}_{\mathrm{ad}} \leftarrow 0$; $\hat{\mathbf{A}}, \hat{\mathbf{S}} \leftarrow \varnothing$ (unset, no identity blend); $\mathbf{A}_{\mathrm{inv}}, \mathbf{S}_{\mathrm{inv}} \leftarrow \mathbf{I}$; $\mathrm{mode} \leftarrow \mathrm{FO}, \tilde{\rho} \leftarrow 1, \tilde{\nu} \leftarrow 0, \mathrm{dwell} \leftarrow 0, \mathrm{lock} \leftarrow 0, \bar{g} \leftarrow 0$; $\mathcal{F}_{\mathrm{collect}} \leftarrow \mathrm{True}; t \leftarrow 0$
\While{not converged}
    \State $t \leftarrow t + 1$; sample minibatch $\mathcal{B}^{(t)}$
    \State $L_{\mathrm{base}} \leftarrow L(\boldsymbol{\theta}^{(t)}; \mathcal{B}^{(t)})$; $\mathbf{g} \leftarrow \nabla_{\boldsymbol{\theta}} L_{\mathrm{base}}$ \Comment{Hooks cache $\mathbf{A}_{\mathrm{raw}}, \mathbf{S}_{\mathrm{raw}}$}
    \If{use\_GC} \State $\mathbf{g}^{(\ell)} \leftarrow \mathrm{CENTRALIZE}(\mathbf{g}^{(\ell)})\ \forall \ell,\ \mathrm{ndim} \ge 2$ \EndIf

    \Statex \textbf{Phase 1: ZO Sensing (every $\tau_{\mathrm{Z}}$ steps)}
    \If{$t \bmod \tau_{\mathrm{Z}} == 0$}
        \State $\mathcal{F}_{\mathrm{collect}} \leftarrow \mathrm{False}$; $S_t \leftarrow$ next $n_p$ blocks (round-robin, all $N$) \Comment{$\hat{\nu}$ gates escape even for $d_\ell > d_{\max}$}
        \For{$\ell \in S_t$}
            \State $(\hat{\rho}, \hat{\nu}) \leftarrow \mathrm{ZO\_PROBE}(\boldsymbol{\theta}^{(t)}, \mathcal{B}^{(t)}, \ell, m, \epsilon_{\mathrm{Z}})$
            \State $\tilde{\rho}^{(\ell)} \leftarrow \alpha_{\mathrm{ema}} \tilde{\rho}^{(\ell)} + (1-\alpha_{\mathrm{ema}}) \hat{\rho}$; $\tilde{\nu}^{(\ell)} \leftarrow \hat{\nu}$
        \EndFor
        \State $\mathcal{F}_{\mathrm{collect}} \leftarrow \mathrm{True}$; $\{\mathrm{mode}^{(\ell)}\} \leftarrow \mathrm{CONTROLLER}(\{\tilde{\rho}\}, \{\mathrm{dwell}\}, \beta_{\mathrm{max}})$
    \EndIf

    \Statex \textbf{Phase 2: ZO Actuator (Saddle Escape)}
    \For{$\ell \in S_t$ sensed at this step}
        \If{$\tilde{\nu}^{(\ell)} \ge \nu_{\mathrm{thr}}$ and $\|\mathbf{g}^{(\ell)}\|_2 \le \gamma_g \bar{g}^{(\ell)}$ and $\mathrm{lock}^{(\ell)} == 0$}
            \State $\mathbf{u} \sim \mathrm{Uniform}(\mathcal{S}^{d_\ell - 1})$; $\boldsymbol{\theta}^{(\ell)} \leftarrow \boldsymbol{\theta}^{(\ell)} + r_{\mathrm{esc}} \|\boldsymbol{\theta}^{(\ell)}\|_2 \mathbf{u}$
            \State $\mathbf{M}^{(\ell)} \leftarrow \frac{1}{2}\mathbf{M}^{(\ell)}$; $\mathbf{M}_{\mathrm{ad}}^{(\ell)} \leftarrow \frac{1}{2}\mathbf{M}_{\mathrm{ad}}^{(\ell)}$; $\mathrm{lock}^{(\ell)} \leftarrow \tau_{\mathrm{esc}}$
        \EndIf
    \EndFor
    \For{all $\ell$} $\mathrm{lock}^{(\ell)} \leftarrow \max(\mathrm{lock}^{(\ell)} - 1, 0)$; $\bar{g}^{(\ell)} \leftarrow \mathrm{EMA}(\|\mathbf{g}^{(\ell)}\|_2)$ \EndFor

    \Statex \textbf{Phase 3: Fisher Factors and Inverses}
    \If{$t \bmod \tau_{\mathrm{stat}} == 0$} \State update $\hat{\mathbf{A}}^{(\ell)}, \hat{\mathbf{S}}^{(\ell)}$ for all eligible $\ell$ via Eq.~\eqref{eq:kfac_ema} \Comment{first observation sets them directly} \EndIf
    \If{$t \bmod \tau_{\mathrm{inv}} == 0$}
        \For{$\ell$ with $\mathrm{mode}^{(\ell)} == \mathrm{SO}$ and $\dim(\ell) \le d_{\mathrm{max}}$}
            \State $\pi \leftarrow \sqrt{(\mathrm{tr}(\hat{\mathbf{A}}) \cdot \dim(\hat{\mathbf{S}})) / (\mathrm{tr}(\hat{\mathbf{S}}) \cdot \dim(\hat{\mathbf{A}}) + \delta)}$
            \State $\mathbf{A}_{\mathrm{inv}}^{(\ell)} \leftarrow (\hat{\mathbf{A}} + \sqrt{\lambda_{\mathrm{damp}}} \pi \mathbf{I})^{-1}$; $\mathbf{S}_{\mathrm{inv}}^{(\ell)} \leftarrow (\hat{\mathbf{S}} + (\sqrt{\lambda_{\mathrm{damp}}}/\pi) \mathbf{I})^{-1}$
        \EndFor
    \EndIf

    \Statex \textbf{Phase 4: Update Engine with Grafting}
    \For{all $\ell$}
        \State $\boldsymbol{\Delta}_{\mathrm{ad}} \leftarrow \textsc{AdamStep}(\mathbf{g}^{(\ell)};\, \mathbf{M}_{\mathrm{ad}}, \mathbf{V}_{\mathrm{ad}}, t)$ \Comment{Eqs.~\eqref{eq:moments}--\eqref{eq:adamw_update}; shadow magnitude}
        \If{$\mathrm{mode}^{(\ell)} == \mathrm{SO}$}
            \State $\mathbf{C} \leftarrow [\mathbf{G}^{(\ell)} \;\; \mathbf{g}_b^{(\ell)}]$; $\mathbf{G}_{\mathrm{pre}}, \mathbf{g}_{b,\mathrm{pre}} \leftarrow \mathrm{Split}(\mathbf{S}_{\mathrm{inv}}^{(\ell)} \mathbf{C} \mathbf{A}_{\mathrm{inv}}^{(\ell)})$
            \State $\boldsymbol{\Delta}_{\mathrm{act}} \leftarrow \textsc{AdamStep}(\mathbf{G}_{\mathrm{pre}};\, \mathbf{M}, \mathbf{V}, t)$ \Comment{separate moments, in the preconditioned basis}
        \Else
            \State $\boldsymbol{\Delta}_{\mathrm{act}} \leftarrow \boldsymbol{\Delta}_{\mathrm{ad}}$; $\mathbf{M} \leftarrow \mathbf{M}_{\mathrm{ad}}$; $\mathbf{V} \leftarrow \mathbf{V}_{\mathrm{ad}}$
        \EndIf
        \State $\boldsymbol{\Delta} \leftarrow \|\boldsymbol{\Delta}_{\mathrm{ad}}\|_F \cdot \boldsymbol{\Delta}_{\mathrm{act}} / (\|\boldsymbol{\Delta}_{\mathrm{act}}\|_F + \delta)$ \Comment{GRAFTING}
        \State $\boldsymbol{\theta}^{(\ell)} \leftarrow \boldsymbol{\theta}^{(\ell)} - \eta^{(t)}(\boldsymbol{\Delta} + \lambda_{\mathrm{wd}}\boldsymbol{\theta}^{(\ell)})$
    \EndFor

    \State \textbf{Non-block params} ($p \notin \{\ell\}$: norm layers, embeddings, biases): plain decoupled AdamW$(\beta_1,\beta_2,\epsilon_a,\lambda_{\mathrm{wd}})$ on $\nabla_p L$
\EndWhile
\end{algorithmic}
\end{algorithm*}

\section{Experimental Setup}
\label{sec:experimental_setup}

To evaluate the optimization capabilities of the proposed Trinity, we evaluated on image classification benchmark across three distinct structural paradigms: (i) a heavy
convolutional neural network (CNN), (ii) a large-scale vision transformer backbone, and (iii) a lightweight residual network.
To ensure the reproducibility of our claims, the complete PyTorch implementation of the Trinity framework, alongside all benchmarking pipelines, hyperparameter search configurations, and ablation scripts, has been open-sourced in an anonymized repository: \url{https://anonymous.4open.science/r/Trinity-opt-DSAA2026/}
\subsection{Dataset and Imbalance Configuration}
For all evaluations, we utilized a challenging Long-Tailed
version of the CIFAR-10 dataset to simulate real-world label imbalances and out-of-distribution tracking performance. We
applied an architectural imbalance factor of $0.01$, creating a severe
$100:1$ ratio between the dominant majority and rare minority classes. This long-tailed pruning significantly restricted
the training capacity to a total of $12{,}406$ active training samples
across all three experimental settings.

\subsection{Architectural Configurations}
\label{subsec:architectures}

To systematically evaluate the scaling properties and geometric adaptability of our proposed framework across highly diverse parameter regimes, we adopt three structural baselines:

\begin{itemize}
    \item \textbf{Experiment 1 (Wide Residual Network):} We evaluate a Wide-ResNet-101-2 architecture containing approximately $126.9\times 10^6$ parameters. Following standard practice for $32 \times 32$ spatial resolutions, the foundational stem of the network was modified to replace the aggressive $7 \times 7$ strided convolution and max-pooling layers with a $3 \times 3$ stride-1 convolution, preserving spatial fidelity.
    
     \item \textbf{Experiment 2 (Vision Transformer):} We configure a custom transformer comprising approximately $10.7\times 10^6$ parameters. The network utilizes a patch tokenization layer executed via a $4\times 4$ convolutional projection (yielding a sequence length of 64 tokens), a learnable cls\_token, and standard positional embeddings. The structural core consists of 6 Transformer encoder blocks operating with an embedding dimension of 384, 6 independent self-attention heads, and an inner feed-forward network (FFN) dimension of 1536.
    
    \item \textbf{Experiment 3 (ResNet-18):} We instantiate a compact ResNet-18 model containing approximately $11\times 10^6$ parameters. This baseline is initialized without pre-trained weights and is mapped to a 10-class output classification head. It serves as a low-complexity control reference to isolate optimizer efficiency across vastly different architectural regimes.
\end{itemize}

All three architectural configurations are evaluated using an identical data pipeline operating with a batch size of $128$. The parameter paths are regularized via a global gradient norm clipping threshold bounded strictly at $1.0$, and the learning rate trajectories are modulated across training epochs using a Cosine Annealing schedule policy.

% \subsection{Hyperparameter Tuning}

% We performed rigorous hyperparameter optimization via the Optuna tuning framework over continuous configuration spaces to isolate peak baseline performance. Across all baseline models, the base learning rates ($\eta$) and weight decay coefficients ($\lambda_{\mathrm{wd}}$) were systematically optimized. For the Trinity architecture, this search domain was expanded to encapsulate its unique multi-order parameters: specifically, the base Tikhonov damping coefficients ($\lambda$) governing second-order K-FAC inversion stability, and the zeroth-order gradient scaling factors ($\zeta$) controlling Random Direction Stochastic Approximation (RDSA) exploration intensity. The discovered optimal hyperparameter sets across the respective architectural spaces are detailed below.
% \smallskip
% \noindent\textbf{Wide-ResNet-101-2:}
% Trinity ($\eta\!\approx\!0.0037$, $\lambda_{wd}\!\approx\!1.4{\times}10^{-4}$,
% damping$\,\approx\!0.0117$, $\tau_{K}\!=\!10$, $\zeta\!\approx\!0.279$)
% ;
% Adam ($\eta\!\approx\!3.4{\times}10^{-4}$,
% $\lambda_{wd}\!\approx\!1.3{\times}10^{-4}$) ;
% RMSProp ($\eta\!\approx\!4.0{\times}10^{-4}$,
% $\lambda_{wd}\!\approx\!1.5{\times}10^{-5}$) ;
% SGD ($\eta\!\approx\!0.006$,
% $\lambda_{wd}\!\approx\!1.18{\times}10^{-4}$) .

% \smallskip
% \noindent\textbf{Vision Transformer:}
% Trinity ($\eta\!\approx\!2.0{\times}10^{-4}$,
% $\lambda_{wd}\!\approx\!2.17{\times}10^{-5}$,
% damping$\,\approx\!0.00137$, $\tau_{K}\!=\!5$, $\zeta\!\approx\!0.299$)
% ;
% Adam ($\eta\!\approx\!4.1{\times}10^{-4}$,
% $\lambda_{wd}\!\approx\!4.6{\times}10^{-4}$) ;
% RMSProp ($\eta\!\approx\!1.1{\times}10^{-4}$,
% $\lambda_{wd}\!\approx\!3.93{\times}10^{-5}$) ;
% SGD ($\eta\!\approx\!0.0075$, $\lambda_{wd}\!\approx\!9.78{\times}10^{-3}$)
% .

% \smallskip
% \noindent\textbf{ResNet-18:}
% Trinity ($\eta\!\approx\!0.00245$,
% $\lambda_{wd}\!\approx\!1.60{\times}10^{-4}$,
% damping$\,\approx\!0.00674$, $\tau_{K}\!=\!5$, $\zeta\!\approx\!0.462$)
% ;
% Adam ($\eta\!\approx\!2.20{\times}10^{-4}$,
% $\lambda_{wd}\!\approx\!5.02{\times}10^{-4}$) ;
% RMSProp ($\eta\!\approx\!2.99{\times}10^{-4}$,
% $\lambda_{wd}\!\approx\!2.02{\times}10^{-5}$) ;
% SGD ($\eta\!\approx\!0.00853$,
% $\lambda_{wd}\!\approx\!5.03{\times}10^{-3}$) .

%------------------------------------------------------------------
\section{Results}
\label{sec:results}

We evaluated our proposed framework across all three deep learning paradigms. Tables~\ref{tab:results_exp1}, \ref{tab:results_exp2}, and \ref{tab:results_exp3} report the peak validation accuracy, final validation accuracy, final training loss, and total wall-clock training time for each optimizer's full 200-epoch training run across Experiment 1 (Wide-ResNet-101-2), Experiment 2 (Vision Transformer), and Experiment 3 (ResNet-18). Figures~\ref{fig:exp1}, \ref{fig:exp2}, and \ref{fig:exp3} plot the corresponding training loss convergence trajectories across all 200 epochs.

\begin{table}[htbp]
\centering
\caption{Performance Comparison on Long-Tailed CIFAR-10: Wide-ResNet-101-2 (Exp.\,1)}
\label{tab:results_exp1}
\setlength{\tabcolsep}{4pt}
\begin{tabular}{lcccc}
\toprule
\textbf{Optimizer} & \textbf{Peak Acc.} & \textbf{Final Acc.} & \textbf{Final Loss} & \textbf{Time\,(s)} \\
\midrule
SGD (Nesterov)      & \textbf{66.17\%} & \textbf{65.02\%} & 0.0025 & 933 \\
RMSProp             & 59.68\% & 58.10\% & 0.0015 & 956 \\
Adam                & 61.79\% & 60.78\% & 0.0008 & 989 \\
\textbf{Trinity}    & 63.48\% & 62.45\% & \textbf{0.0006} & 13263 \\
\bottomrule
\end{tabular}
\end{table}

\begin{table}[htbp]
\centering
\caption{Performance Comparison on Long-Tailed CIFAR-10: Vision Transformer (Exp.\,2)}
\label{tab:results_exp2}
\setlength{\tabcolsep}{4pt}
\begin{tabular}{lcccc}
\toprule
\textbf{Optimizer} & \textbf{Peak Acc.} & \textbf{Final Acc.} & \textbf{Final Loss} & \textbf{Time\,(s)} \\
\midrule
SGD (Nesterov)      & 51.22\% & \textbf{50.59\%} & 0.2922 & 544 \\
RMSProp             & 46.95\% & 43.08\% & 0.0646 & 555 \\
Adam                & \textbf{52.70\%} & 48.37\% & 0.0056 & 555 \\
\textbf{Trinity}    & 48.86\% & 47.37\% & \textbf{0.0001} & 1396 \\
\bottomrule
\end{tabular}
\end{table}

\begin{table}[htbp]
\centering
\caption{Performance Comparison on Long-Tailed CIFAR-10: ResNet-18 (Exp.\,3)}
\label{tab:results_exp3}
\setlength{\tabcolsep}{4pt}
\begin{tabular}{lcccc}
\toprule
\textbf{Optimizer} & \textbf{Peak Acc.} & \textbf{Final Acc.} & \textbf{Final Loss} & \textbf{Time\,(s)} \\
\midrule
SGD (Nesterov)      & \textbf{63.82\%} & \textbf{60.86\%} & 0.0027 & 435 \\
RMSProp             & 61.46\% & 59.62\% & 0.0026 & 427 \\
Adam                & 62.25\% & 59.56\% & 0.0017 & 395 \\
\textbf{Trinity}    & 61.17\% & 59.85\% & \textbf{0.0011} & 1026 \\
\bottomrule
\end{tabular}
\end{table}

The empirical findings demonstrate that Trinity consistently establishes a superior optimization capability, achieving the lowest training loss across all architectures, while exhibiting nuanced structural trade-offs in validation accuracy between parameter-dense and compact models.

In \textbf{Experiment 1} (Figure~\ref{fig:exp1} and Table~\ref{tab:results_exp1}), training the massive Wide-ResNet-101-2 architecture from scratch on severely imbalanced data demonstrates that Trinity achieves the lowest final training loss of $\mathbf{0.0006}$ (with an absolute minimum of $0.0002$), significantly lower than Adam ($0.0008$), RMSProp ($0.0015$), and SGD ($0.0025$). In terms of validation generalization, SGD achieves the highest peak accuracy of $66.17\%$, while Trinity reaches a highly competitive peak accuracy of $63.48\%$ and final accuracy of $62.45\%$, noticeably outperforming both Adam ($61.79\%$) and RMSProp ($59.68\%$). 

In \textbf{Experiment 2} (Figure~\ref{fig:exp2} and Table~\ref{tab:results_exp2}), the non-local attention mechanics and sequence space of the Vision Transformer highlight Trinity's distinct optimization power. Trinity drives training loss to an empirical floor of $\mathbf{0.0001}$, whereas standard first-order methods encounter severe landscape ill-conditioning under the 100:1 label imbalance (SGD stalls at $0.2922$, RMSProp settles at $0.0646$, and Adam reaches $0.0056$). Trinity minimizes the training objective up to three orders of magnitude deeper than first-order baselines. In validation accuracy over 200 epochs, Adam reaches a peak of $52.70\%$, SGD achieves $51.22\%$, Trinity attains $48.86\%$ ($47.37\%$ final), and RMSProp reaches $46.95\%$. The periodic calculation of K-FAC factor inversions and auxiliary RDSA forward evaluations incurs a deterministic wall-clock overhead ($1396\text{s}$ vs. $555\text{s}$ for Adam), but enables Trinity to overcome sharp ill-conditioned barriers that hinder traditional optimizers.

In \textbf{Experiment 3} (Figure~\ref{fig:exp3} and Table~\ref{tab:results_exp3}), the compact ResNet-18 backbone exhibits a different dynamic. Trinity again establishes the absolute lowest training loss at $\mathbf{0.0011}$ (minimum $0.0009$), outperforming Adam ($0.0017$), RMSProp ($0.0026$), and SGD ($0.0027$). In validation accuracy, SGD reaches $63.82\%$, Adam attains $62.25\%$, RMSProp reaches $61.46\%$, and Trinity achieves $61.17\%$ peak ($59.85\%$ final, surpassing Adam's final $59.56\%$).

\definecolor{colorhybrid}{RGB}{31, 119, 180}   % Modern Steel Blue
\definecolor{coloradam}{RGB}{214, 39, 40}     % Academic Muted Red
\definecolor{colorrmsprop}{RGB}{44, 160, 44}  % Balanced Forest Green
\definecolor{colorsgd}{RGB}{255, 127, 14}     % Safety Orange

\pgfplotsset{compat=1.18}

% ──────────────────────────────────────────────────────────────────────────────
% FIGURE 1: Experiment 1 (Wide-ResNet-101-2) Training Loss
% ──────────────────────────────────────────────────────────────────────────────
\begin{figure*}[!t]
  \centering
  \begin{tikzpicture}
    % --- LOSS PLOT ---
    \begin{axis}[
      name=lossplot,
      width=0.48\textwidth,
      height=5.0cm,
      ymode=log,
      xlabel={Epoch},
      ylabel={Training Loss},
      xlabel style={font=\small, yshift=2pt},
      ylabel style={font=\small, xshift=-2pt},
      tick label style={font=\footnotesize},
      xmin=1,   xmax=200,
      ymin=0.0001, ymax=4.0,
      xtick={1,25,50,75,100,125,150,175,200},
      minor tick num=4,
      grid=none,
      legend pos=south west,
      legend style={
        font=\footnotesize,
        fill=white,
        fill opacity=0.95,
        text opacity=1,
        row sep=2pt,
        inner sep=4pt,
        rounded corners=1pt
      },
      clip=true,
    ]
    \addplot[line width=1.3pt, color=colorhybrid, mark=none] table [x=epoch, y=loss, col sep=comma] {CSVs_clean/exp1_trinity.csv};
    \addplot[line width=1.3pt, color=coloradam, mark=none] table [x=epoch, y=loss, col sep=comma] {CSVs_clean/exp1_adam.csv};
    \addplot[line width=1.3pt, color=colorrmsprop, mark=none] table [x=epoch, y=loss, col sep=comma] {CSVs_clean/exp1_rmsprop.csv};
    \addplot[line width=1.3pt, color=colorsgd, mark=none] table [x=epoch, y=loss, col sep=comma] {CSVs_clean/exp1_sgd.csv};
    
    \legend{Trinity (Proposed), Adam, RMSProp, SGD}
    \end{axis}

    % --- ACCURACY PLOT ---
    \begin{axis}[
      at={(lossplot.east)},
      anchor=west,
      xshift=1.0cm,
      width=0.48\textwidth,
      height=5.0cm,
      xlabel={Epoch},
      ylabel={Validation Accuracy (\%)},
      xlabel style={font=\small, yshift=2pt},
      ylabel style={font=\small, xshift=-2pt},
      tick label style={font=\footnotesize},
      xmin=1,   xmax=200,
      ymin=10, ymax=70,
      xtick={1,25,50,75,100,125,150,175,200},
      grid=both,
      legend pos=south east,
      legend style={
        font=\footnotesize,
        fill=white,
        fill opacity=0.95,
        text opacity=1,
        row sep=2pt,
        inner sep=4pt,
        rounded corners=1pt
      },
      clip=true,
    ]
    \addplot[line width=1.3pt, color=colorhybrid, mark=none] table [x=epoch, y=val_acc, col sep=comma] {CSVs_clean/exp1_trinity.csv};
    \addplot[line width=1.3pt, color=coloradam, mark=none] table [x=epoch, y=val_acc, col sep=comma] {CSVs_clean/exp1_adam.csv};
    \addplot[line width=1.3pt, color=colorrmsprop, mark=none] table [x=epoch, y=val_acc, col sep=comma] {CSVs_clean/exp1_rmsprop.csv};
    \addplot[line width=1.3pt, color=colorsgd, mark=none] table [x=epoch, y=val_acc, col sep=comma] {CSVs_clean/exp1_sgd.csv};
    \end{axis}
  \end{tikzpicture}
  \caption{Experiment 1 (Wide-ResNet-101-2) convergence: Training Loss (left) and Validation Accuracy (right).}
  \label{fig:exp1}
\end{figure*}

% ──────────────────────────────────────────────────────────────────────────────
% FIGURE 2: Experiment 2 (Vision Transformer) Training Loss
% ──────────────────────────────────────────────────────────────────────────────
\begin{figure*}[!t]
  \centering
  \begin{tikzpicture}
    % --- LOSS PLOT ---
    \begin{axis}[
      name=lossplot,
      width=0.48\textwidth,
      height=5.0cm,
      ymode=log,
      xlabel={Epoch},
      ylabel={Training Loss},
      xlabel style={font=\small, yshift=2pt},
      ylabel style={font=\small, xshift=-2pt},
      tick label style={font=\footnotesize},
      xmin=1,   xmax=200,
      ymin=0.00005, ymax=2.5,
      xtick={1,25,50,75,100,125,150,175,200},
      minor tick num=4,
      grid=none,
      legend pos=south west,
      legend style={
        font=\footnotesize,
        fill=white,
        fill opacity=0.95,
        text opacity=1,
        row sep=2pt,
        inner sep=4pt,
        rounded corners=1pt
      },
      clip=true,
    ]
    \addplot[line width=1.3pt, color=colorhybrid, mark=none] table [x=epoch, y=loss, col sep=comma] {CSVs_clean/exp2_trinity.csv};
    \addplot[line width=1.3pt, color=coloradam, mark=none] table [x=epoch, y=loss, col sep=comma] {CSVs_clean/exp2_adam.csv};
    \addplot[line width=1.3pt, color=colorrmsprop, mark=none] table [x=epoch, y=loss, col sep=comma] {CSVs_clean/exp2_rmsprop.csv};
    \addplot[line width=1.3pt, color=colorsgd, mark=none] table [x=epoch, y=loss, col sep=comma] {CSVs_clean/exp2_sgd.csv};
    
    \legend{Trinity (Proposed), Adam, RMSProp, SGD}
    \end{axis}

    % --- ACCURACY PLOT ---
    \begin{axis}[
      at={(lossplot.east)},
      anchor=west,
      xshift=1.0cm,
      width=0.48\textwidth,
      height=5.0cm,
      xlabel={Epoch},
      ylabel={Validation Accuracy (\%)},
      xlabel style={font=\small, yshift=2pt},
      ylabel style={font=\small, xshift=-2pt},
      tick label style={font=\footnotesize},
      xmin=1,   xmax=200,
      ymin=10, ymax=60,
      xtick={1,25,50,75,100,125,150,175,200},
      grid=both,
      legend pos=south east,
      legend style={
        font=\footnotesize,
        fill=white,
        fill opacity=0.95,
        text opacity=1,
        row sep=2pt,
        inner sep=4pt,
        rounded corners=1pt
      },
      clip=true,
    ]
    \addplot[line width=1.3pt, color=colorhybrid, mark=none] table [x=epoch, y=val_acc, col sep=comma] {CSVs_clean/exp2_trinity.csv};
    \addplot[line width=1.3pt, color=coloradam, mark=none] table [x=epoch, y=val_acc, col sep=comma] {CSVs_clean/exp2_adam.csv};
    \addplot[line width=1.3pt, color=colorrmsprop, mark=none] table [x=epoch, y=val_acc, col sep=comma] {CSVs_clean/exp2_rmsprop.csv};
    \addplot[line width=1.3pt, color=colorsgd, mark=none] table [x=epoch, y=val_acc, col sep=comma] {CSVs_clean/exp2_sgd.csv};
    \end{axis}
  \end{tikzpicture}
  \caption{Experiment 2 (Vision Transformer) convergence: Training Loss (left) and Validation Accuracy (right).}
  \label{fig:exp2}
\end{figure*}

% ──────────────────────────────────────────────────────────────────────────────
% FIGURE 3: Experiment 3 (ResNet-18) Training Loss
% ──────────────────────────────────────────────────────────────────────────────
\begin{figure*}[!t]
  \centering
  \begin{tikzpicture}
    % --- LOSS PLOT ---
    \begin{axis}[
      name=lossplot,
      width=0.48\textwidth,
      height=5.0cm,
      ymode=log,
      xlabel={Epoch},
      ylabel={Training Loss},
      xlabel style={font=\small, yshift=2pt},
      ylabel style={font=\small, xshift=-2pt},
      tick label style={font=\footnotesize},
      xmin=1,   xmax=200,
      ymin=0.0005, ymax=3.0,
      xtick={1,25,50,75,100,125,150,175,200},
      minor tick num=4,
      grid=none,
      legend pos=south west,
      legend style={
        font=\footnotesize,
        fill=white,
        fill opacity=0.95,
        text opacity=1,
        row sep=2pt,
        inner sep=4pt,
        rounded corners=1pt
      },
      clip=true,
    ]
    \addplot[line width=1.3pt, color=colorhybrid, mark=none] table [x=epoch, y=loss, col sep=comma] {CSVs_clean/exp3_trinity.csv};
    \addplot[line width=1.3pt, color=coloradam, mark=none] table [x=epoch, y=loss, col sep=comma] {CSVs_clean/exp3_adam.csv};
    \addplot[line width=1.3pt, color=colorrmsprop, mark=none] table [x=epoch, y=loss, col sep=comma] {CSVs_clean/exp3_rmsprop.csv};
    \addplot[line width=1.3pt, color=colorsgd, mark=none] table [x=epoch, y=loss, col sep=comma] {CSVs_clean/exp3_sgd.csv};
    
    \legend{Trinity (Proposed), Adam, RMSProp, SGD}
    \end{axis}

    % --- ACCURACY PLOT ---
    \begin{axis}[
      at={(lossplot.east)},
      anchor=west,
      xshift=1.0cm,
      width=0.48\textwidth,
      height=5.0cm,
      xlabel={Epoch},
      ylabel={Validation Accuracy (\%)},
      xlabel style={font=\small, yshift=2pt},
      ylabel style={font=\small, xshift=-2pt},
      tick label style={font=\footnotesize},
      xmin=1,   xmax=200,
      ymin=10, ymax=70,
      xtick={1,25,50,75,100,125,150,175,200},
      grid=both,
      legend pos=south east,
      legend style={
        font=\footnotesize,
        fill=white,
        fill opacity=0.95,
        text opacity=1,
        row sep=2pt,
        inner sep=4pt,
        rounded corners=1pt
      },
      clip=true,
    ]
    \addplot[line width=1.3pt, color=colorhybrid, mark=none] table [x=epoch, y=val_acc, col sep=comma] {CSVs_clean/exp3_trinity.csv};
    \addplot[line width=1.3pt, color=coloradam, mark=none] table [x=epoch, y=val_acc, col sep=comma] {CSVs_clean/exp3_adam.csv};
    \addplot[line width=1.3pt, color=colorrmsprop, mark=none] table [x=epoch, y=val_acc, col sep=comma] {CSVs_clean/exp3_rmsprop.csv};
    \addplot[line width=1.3pt, color=colorsgd, mark=none] table [x=epoch, y=val_acc, col sep=comma] {CSVs_clean/exp3_sgd.csv};
    \end{axis}
  \end{tikzpicture}
  \caption{Experiment 3 (ResNet-18) convergence: Training Loss (left) and Validation Accuracy (right).}
  \label{fig:exp3}
\end{figure*}

%-----------------------------------------------------
\section{Analysis and Discussion}
\label{sec:analysis}


To systematically evaluate the individual contributions of Trinity's algorithmic components, we conducted extensive ablation experiments across all three architectural paradigms evaluated in this work: a high-capacity deep convolutional network (\textbf{Wide-ResNet-101-2}, Exp.\,1), a multi-head self-attention architecture (\textbf{Vision Transformer}, Exp.\,2), and a residual network (\textbf{ResNet-18}, Exp.\,3). All models were trained from scratch over a full \textbf{200-epoch} budget on Long-Tailed CIFAR-10 under severe class imbalance ($\text{IF}=0.01$, 100:1 head-to-tail imbalance ratio, $12{,}406$ training images) using a Cosine Annealing learning rate schedule and batch size 128.

We evaluated seven systematically controlled ablation arms to isolate each modular mechanism:
\begin{enumerate}
    \item \textbf{Adam Baseline}: Standard first-order adaptive momentum optimizer.
    \item \textbf{Trinity (Full)}: Complete framework with zeroth-order condition monitoring, selective K-FAC curvature preconditioning, Adam magnitude grafting, gradient centralization, and saddle escape tracking.
    \item \textbf{Trinity w/o K-FAC}: Disabling second-order curvature preconditioning ($\text{SO}=0$), restricting updates to first-order adaptive momentum combined with zeroth-order condition monitoring and gradient centralization.
    \item \textbf{Always-On K-FAC}: Forcing second-order preconditioning across all network layers at every epoch, bypassing Trinity's selective zeroth-order gating.
    \item \textbf{Trinity w/o Magnitude Grafting}: Omitting step-norm projection onto Adam's momentum scale, executing raw second-order natural gradient updates on preconditioned layers.
    \item \textbf{Trinity w/o Gradient Centralization}: Disabling weight vector mean-centering projections.
    \item \textbf{Trinity w/o Saddle Escape}: Disabling zeroth-order exploratory subspace perturbations.
\end{enumerate}

Table~\ref{tab:ablation_comparison} reports the empirical validation ceiling, final validation accuracy at epoch 200, and final training loss across all three architectural paradigms for all seven ablation configurations.

\begin{table}[htbp]
\centering
\caption{Systematic Component Ablation Study Across 200 Epochs (CIFAR-10 LT $\text{IF}=0.01$). Format: \textbf{Peak Val.\% / Final Val.\% / Final Loss}}
\label{tab:ablation_comparison}
\setlength{\tabcolsep}{2.5pt}
\resizebox{\columnwidth}{!}{%
\begin{tabular}{lccc}
\toprule
\textbf{Ablation Variant} & \textbf{Exp 1 (WResNet)} & \textbf{Exp 2 (Transformer)} & \textbf{Exp 3 (ResNet-18)} \\
\midrule
Adam (Baseline)           & 61.65 / 60.07 / 0.0015 & 52.52 / 48.47 / 0.0062 & 61.95 / 59.92 / 0.0023 \\
\textbf{Trinity (Full)}   & \textbf{62.84 / 61.24 / 0.0010} & \textbf{48.93 / 47.27 / 0.0001} & \textbf{61.09 / 59.19 / 0.0018} \\
w/o K-FAC (SO off)     & 62.78 / 61.57 / 0.0006 & 50.00 / 47.43 / 0.0001 & 62.63 / 61.14 / 0.0016 \\
Always-On K-FAC        & 64.64 / 63.13 / 0.0004 & 49.08 / 46.66 / 0.0004 & 62.75 / 61.92 / 0.0018 \\
w/o Magnitude Grafting & 63.25 / 62.77 / 0.0007 & 47.18 / 45.52 / 0.0010 & 61.64 / 60.13 / 0.0035 \\
w/o Grad. Centralization & 62.58 / 60.30 / 0.0002 & 49.27 / 47.21 / 0.0001 & 60.62 / 58.64 / 0.0008 \\
w/o Saddle Escape      & 63.74 / 62.36 / 0.0007 & 48.85 / 46.72 / 0.0002 & 60.55 / 59.30 / 0.0018 \\
\bottomrule
\end{tabular}%
}
\end{table}

\subsection{Component-Wise Architectural Insights}

\textbf{Step Magnitude Grafting as a Non-Negotiable Stabilizer.}
A paramount finding of our ablation suite is the use of step magnitude grafting. When second-order natural gradient steps are computed for K-FAC layers, their unconstrained Euclidean norms diverge from the adaptive momentum updates of first-order parameters. Omitting grafting (\textit{Trinity w/o Magnitude Grafting}) can affect the performance. For Vision transformer the final loss magnitude at $0.0010$ (versus $0.0001$). Consequently, its peak ($47.18\%$) and final ($45.52\%$) validation accuracies are the lowest among all evaluated Trinity variants. On ResNet-18, disabling grafting similarly stalls late-stage convergence, yielding the highest final loss ($0.0035$) of any Trinity configuration. Grafting functions as a vital metric stabilizer that prevents gradient scale dissonance across heterogeneous layers.

\textbf{Selective Curvature Gating vs.\ Always-On K-FAC.}
In standard second-order methods, K-FAC inversions are performed indiscriminately across all layers, incurring steep computational costs. In Wide-ResNet-101-2 (105 parameter blocks), \textit{Always-On K-FAC} requires $166.8$ minutes ($2.78$\,hours) due to inverting massive Kronecker factors every epoch. In contrast, Trinity's selective zeroth-order gating makes high-capacity training computationally practical ($37.3$ minutes) while attaining strong final accuracy ($61.24\%$). In the Vision Transformer benchmark, \textit{Always-On K-FAC} requires $24.6$ minutes ($7.37$\,s/epoch). Trinity's zeroth-order condition sensor continuously evaluates block-level progress metrics ($\rho$), selectively engaging curvature preconditioning on only the 8 most ill-conditioned blocks. This dynamic gating reduces training wall-clock time by $11\%$ ($22.1$ minutes, $6.64$\,s/epoch) while simultaneously improving final validation generalization from $46.66\%$ to $47.27\%$, demonstrating that selective gating actively mitigates the risk of over-preconditioning noisy attention projections. On ResNet-18, selective K-FAC slashes total training time by $40\%$ compared to Always-On K-FAC ($18.6$ minutes vs.\ $31.0$ minutes; $5.57$\,s vs.\ $9.30$\,s/epoch), while maintaining competitive generalization ($59.19\%$ vs.\ $61.92\%$).

\textbf{Architectural Divergence of Gradient Centralization.}
Our ablation exposes an instructive contrast between convolutional and transformer architectures regarding Gradient Centralization (\texttt{use\_gc}). On both convolutional networks, removing GC (\textit{w/o Grad.\ Centralization}) produces the lowest final validation accuracy among all ablation arms: dropping to $58.64\%$ on ResNet-18 and $60.30\%$ on Wide-ResNet-101-2. This confirms that zero-centering gradient vectors along weight vectors provides essential manifold regularization for 2D convolutional filter banks. Conversely, for the Vision Transformer, disabling GC incurs no generalization degradation ($47.21\%$ final accuracy, $49.27\%$ peak) while reducing per-epoch latency from $6.64$\,s to $6.10$\,s ($\approx 8\%$ throughput speedup). This confirms that GC's inductive bias is specific to spatial convolutional kernels and is redundant for linear projection matrices in self-attention.

\textbf{Saddle Escape and Long-Tailed Class Representation.}
Disabling the zeroth-order saddle escape mechanism (\textit{w/o Saddle Escape}) leads to measurable performance degradations across all architectures ($62.36\%$ on Wide-ResNet; $59.30\%$ on ResNet-18; $46.72\%$ on ViT). On the Vision Transformer, class-group breakdown reveals that omitting escape reduces accuracy from $47.27\%$ down to $46.72\%$.


% \section{Conclusion}
% \label{sec:conclusion}
\section{Conclusion}
\label{sec:conclusion}

In this paper, we introduced \textbf{Trinity}, an optimization framework that unifies zeroth-, first-, and second-order mechanics through a principled division of labor: a lightweight forward-only zeroth-order sensor measures Fisher anisotropy at minimal overhead, selectively triggering second-order K-FAC inversions only on severely ill-conditioned parameter blocks while an AdamW base engine guarantees scale-safe updates via magnitude grafting and zeroth-order saddle escape. Evaluated across Wide-ResNet-101-2, Vision Transformers, and ResNet-18 under extreme 100:1 long-tailed label imbalance, Trinity consistently achieves the deepest training loss minimization across all architectures while delivering competitive validation generalization. By proving that geometry-aware preconditioning can be dynamically steered via forward-only runtime sensing rather than applied uniformly, Trinity provides a practical, scalable path for higher-order optimization in modern deep neural networks.


\bibliographystyle{plain}
\begin{thebibliography}{12}

\bibitem{kingma2014adam}
D.~P. Kingma and J.~Ba, ``Adam: A method for stochastic optimization,'' in \emph{International Conference on Learning Representations (ICLR)}, 2015.

\bibitem{loshchilov2019decoupled}
I.~Loshchilov and F.~Hutter, ``Decoupled weight decay regularization,'' in \emph{International Conference on Learning Representations (ICLR)}, 2019.

\bibitem{yong2020gradient}
H.~Yong, J.~Huang, X.~Hua, and L.~Zhang, ``Gradient centralization: A new optimization technique for deep neural networks,'' in \emph{European Conference on Computer Vision (ECCV)}, 2020, pages 635--652.

\bibitem{amari1998natural}
S.-I. Amari, ``Natural gradient works efficiently in learning,'' \emph{Neural Computation}, vol. 10, no. 2, pages 251--276, 1998.

\bibitem{martens2015optimizing}
J.~Martens and R.~Grosse, ``Optimizing neural networks with Kronecker-factored approximate curvature,'' in \emph{International Conference on Machine Learning (ICML)}, 2015, pages 2408--2417.

\bibitem{osawa2019large}
K.~Osawa, Y.~Tsuji, Y.~Ueno, A.~Naruse, R.~Yokota, and S.~Matsuoka, ``Large-scale distributed second-order optimization using Kronecker-factored approximate curvature for deep convolutional neural networks,'' in \emph{Proceedings of the IEEE/CVF Conference on Computer Vision and Pattern Recognition (CVPR)}, 2019.

\bibitem{spall1992multivariate}
J.~C. Spall, ``Multivariate stochastic approximation using a simultaneous perturbation gradient approximation,'' \emph{IEEE Transactions on Automatic Control}, vol. 37, no. 3, pages 332--341, 1992.

\bibitem{nesterov2017random}
Y.~Nesterov and V.~Spokoiny, ``Random gradient-free minimization of convex functions,'' \emph{Foundations of Computational Mathematics}, vol. 17, no. 2, pp. 527--566, 2017.
% \bibitem{foret2020sharpness}
% P.~Foret, A.~Kleiner, H.~Mobahi, and B.~Neyshabur, ``Sharpness-aware minimization for efficiently improving generalization,'' in \emph{International Conference on Learning Representations (ICLR)}, 2021.

% \bibitem{kwon2021asam}
% J.~Kwon, J.~Kim, H.~Park, and I.~K. Choi, ``ASAM: Adaptive sharpness-aware minimization for scale-invariant learning of deep neural networks,'' in \emph{International Conference on Machine Learning (ICML)}, 2021, pages 5905--5914.

\bibitem{zhang2019lookahead}
M.~Zhang, J.~Lucas, J.~Ba, and G.~E. Hinton, ``Lookahead optimizer: k steps forward, 1 step back,'' in \emph{Advances in Neural Information Processing Systems (NeurIPS)}, 2019.

\bibitem{sst2}
R.~Socher, A.~Perelygin, J.~Wu, J.~Chuang, C.~D. Manning, A.~Ng, and C.~Potts, ``Recursive deep models for semantic compositionality over a sentiment treebank,'' in \emph{Proceedings of the Conference on Empirical Methods in Natural Language Processing (EMNLP)}, 2013.

\bibitem{kunstner2024heavy}
F.~Kunstner, R.~Yadav, A.~Milligan, M.~Schmidt, and A.~Bietti, ``Heavy-Tailed Class Imbalance and Why Adam Outperforms Gradient Descent on Language Models,'' \emph{arXiv preprint arXiv:2402.19449}, 2024.


\end{thebibliography}









% \section{Introduction}
% This document is a model and instructions for \LaTeX.
% Please observe the conference page limits. 

% \section{Ease of Use}

% \subsection{Maintaining the Integrity of the Specifications}

% The IEEEtran class file is used to format your paper and style the text. All margins, 
% column widths, line spaces, and text fonts are prescribed; please do not 
% alter them. You may note peculiarities. For example, the head margin
% measures proportionately more than is customary. This measurement 
% and others are deliberate, using specifications that anticipate your paper 
% as one part of the entire proceedings, and not as an independent document. 
% Please do not revise any of the current designations.

% \section{Prepare Your Paper Before Styling}
% Before you begin to format your paper, first write and save the content as a 
% separate text file. Complete all content and organizational editing before 
% formatting. Please note sections \ref{AA}--\ref{SCM} below for more information on 
% proofreading, spelling and grammar.

% Keep your text and graphic files separate until after the text has been 
% formatted and styled. Do not number text heads---{\LaTeX} will do that 
% for you.

% \subsection{Abbreviations and Acronyms}\label{AA}
% Define abbreviations and acronyms the first time they are used in the text, 
% even after they have been defined in the abstract. Abbreviations such as 
% IEEE, SI, MKS, CGS, ac, dc, and rms do not have to be defined. Do not use 
% abbreviations in the title or heads unless they are unavoidable.

% \subsection{Units}
% \begin{itemize}
% \item Use either SI (MKS) or CGS as primary units. (SI units are encouraged.) English units may be used as secondary units (in parentheses). An exception would be the use of English units as identifiers in trade, such as ``3.5-inch disk drive''.
% \item Avoid combining SI and CGS units, such as current in amperes and magnetic field in oersteds. This often leads to confusion because equations do not balance dimensionally. If you must use mixed units, clearly state the units for each quantity that you use in an equation.
% \item Do not mix complete spellings and abbreviations of units: ``Wb/m\textsuperscript{2}'' or ``webers per square meter'', not ``webers/m\textsuperscript{2}''. Spell out units when they appear in text: ``. . . a few henries'', not ``. . . a few H''.
% \item Use a zero before decimal points: ``0.25'', not ``.25''. Use ``cm\textsuperscript{3}'', not ``cc''.)
% \end{itemize}

% \subsection{Equations}
% Number equations consecutively. To make your 
% equations more compact, you may use the solidus (~/~), the exp function, or 
% appropriate exponents. Italicize Roman symbols for quantities and variables, 
% but not Greek symbols. Use a long dash rather than a hyphen for a minus 
% sign. Punctuate equations with commas or periods when they are part of a 
% sentence, as in:
% \begin{equation}
% a+b=\gamma\label{eq}
% \end{equation}

% Be sure that the 
% symbols in your equation have been defined before or immediately following 
% the equation. Use ``\eqref{eq}'', not ``Eq.~\eqref{eq}'' or ``equation \eqref{eq}'', except at 
% the beginning of a sentence: ``Equation \eqref{eq} is . . .''

% \subsection{\LaTeX-Specific Advice}

% Please use ``soft'' (e.g., \verb|\eqref{Eq}|) cross references instead
% of ``hard'' references (e.g., \verb|(1)|). That will make it possible
% to combine sections, add equations, or change the order of figures or
% citations without having to go through the file line by line.

% Please don't use the \verb|{eqnarray}| equation environment. Use
% \verb|{align}| or \verb|{IEEEeqnarray}| instead. The \verb|{eqnarray}|
% environment leaves unsightly spaces around relation symbols.

% Please note that the \verb|{subequations}| environment in {\LaTeX}
% will increment the main equation counter even when there are no
% equation numbers displayed. If you forget that, you might write an
% article in which the equation numbers skip from (17) to (20), causing
% the copy editors to wonder if you've discovered a new method of
% counting.

% {\BibTeX} does not work by magic. It doesn't get the bibliographic
% data from thin air but from .bib files. If you use {\BibTeX} to produce a
% bibliography you must send the .bib files. 

% {\LaTeX} can't read your mind. If you assign the same label to a
% subsubsection and a table, you might find that Table I has been cross
% referenced as Table IV-B3. 

% {\LaTeX} does not have precognitive abilities. If you put a
% \verb|\label| command before the command that updates the counter it's
% supposed to be using, the label will pick up the last counter to be
% cross referenced instead. In particular, a \verb|\label| command
% should not go before the caption of a figure or a table.

% Do not use \verb|\nonumber| inside the \verb|{array}| environment. It
% will not stop equation numbers inside \verb|{array}| (there won't be
% any anyway) and it might stop a wanted equation number in the
% surrounding equation.

% \subsection{Some Common Mistakes}\label{SCM}
% \begin{itemize}
% \item The word ``data'' is plural, not singular.
% \item The subscript for the permeability of vacuum $\mu_{0}$, and other common scientific constants, is zero with subscript formatting, not a lowercase letter ``o''.
% \item In American English, commas, semicolons, periods, question and exclamation marks are located within quotation marks only when a complete thought or name is cited, such as a title or full quotation. When quotation marks are used, instead of a bold or italic typeface, to highlight a word or phrase, punctuation should appear outside of the quotation marks. A parenthetical phrase or statement at the end of a sentence is punctuated outside of the closing parenthesis (like this). (A parenthetical sentence is punctuated within the parentheses.)
% \item A graph within a graph is an ``inset'', not an ``insert''. The word alternatively is preferred to the word ``alternately'' (unless you really mean something that alternates).
% \item Do not use the word ``essentially'' to mean ``approximately'' or ``effectively''.
% \item In your paper title, if the words ``that uses'' can accurately replace the word ``using'', capitalize the ``u''; if not, keep using lower-cased.
% \item Be aware of the different meanings of the homophones ``affect'' and ``effect'', ``complement'' and ``compliment'', ``discreet'' and ``discrete'', ``principal'' and ``principle''.
% \item Do not confuse ``imply'' and ``infer''.
% \item The prefix ``non'' is not a word; it should be joined to the word it modifies, usually without a hyphen.
% \item There is no period after the ``et'' in the Latin abbreviation ``et al.''.
% \item The abbreviation ``i.e.'' means ``that is'', and the abbreviation ``e.g.'' means ``for example''.
% \end{itemize}
% An excellent style manual for science writers is \cite{b7}.

% \subsection{Authors and Affiliations}
% \textbf{The class file is designed for, but not limited to, six authors.} A 
% minimum of one author is required for all conference articles. Author names 
% should be listed starting from left to right and then moving down to the 
% next line. This is the author sequence that will be used in future citations 
% and by indexing services. Names should not be listed in columns nor group by 
% affiliation. Please keep your affiliations as succinct as possible (for 
% example, do not differentiate among departments of the same organization).

% \subsection{Identify the Headings}
% Headings, or heads, are organizational devices that guide the reader through 
% your paper. There are two types: component heads and text heads.

% Component heads identify the different components of your paper and are not 
% topically subordinate to each other. Examples include Acknowledgments and 
% References and, for these, the correct style to use is ``Heading 5''. Use 
% ``figure caption'' for your Figure captions, and ``table head'' for your 
% table title. Run-in heads, such as ``Abstract'', will require you to apply a 
% style (in this case, italic) in addition to the style provided by the drop 
% down menu to differentiate the head from the text.

% Text heads organize the topics on a relational, hierarchical basis. For 
% example, the paper title is the primary text head because all subsequent 
% material relates and elaborates on this one topic. If there are two or more 
% sub-topics, the next level head (uppercase Roman numerals) should be used 
% and, conversely, if there are not at least two sub-topics, then no subheads 
% should be introduced.

% \subsection{Figures and Tables}
% \paragraph{Positioning Figures and Tables} Place figures and tables at the top and 
% bottom of columns. Avoid placing them in the middle of columns. Large 
% figures and tables may span across both columns. Figure captions should be 
% below the figures; table heads should appear above the tables. Insert 
% figures and tables after they are cited in the text. Use the abbreviation 
% ``Fig.~\ref{fig}'', even at the beginning of a sentence.

% \begin{table}[htbp]
% \caption{Table Type Styles}
% \begin{center}
% \begin{tabular}{|c|c|c|c|}
% \hline
% \textbf{Table}&\multicolumn{3}{|c|}{\textbf{Table Column Head}} \\
% \cline{2-4} 
% \textbf{Head} & \textbf{\textit{Table column subhead}}& \textbf{\textit{Subhead}}& \textbf{\textit{Subhead}} \\
% \hline
% copy& More table copy$^{\mathrm{a}}$& &  \\
% \hline
% \multicolumn{4}{l}{$^{\mathrm{a}}$Sample of a Table footnote.}
% \end{tabular}
% \label{tab1}
% \end{center}
% \end{table}

% \begin{figure}[htbp]
% \centerline{\includegraphics{fig1.png}}
% \caption{Example of a figure caption.}
% \label{fig}
% \end{figure}

% Figure Labels: Use 8 point Times New Roman for Figure labels. Use words 
% rather than symbols or abbreviations when writing Figure axis labels to 
% avoid confusing the reader. As an example, write the quantity 
% ``Magnetization'', or ``Magnetization, M'', not just ``M''. If including 
% units in the label, present them within parentheses. Do not label axes only 
% with units. In the example, write ``Magnetization (A/m)'' or ``Magnetization 
% \{A[m(1)]\}'', not just ``A/m''. Do not label axes with a ratio of 
% quantities and units. For example, write ``Temperature (K)'', not 
% ``Temperature/K''.

% \section*{Acknowledgment}

% The preferred spelling of the word ``acknowledgment'' in America is without 
% an ``e'' after the ``g''. Avoid the stilted expression ``one of us (R. B. 
% G.) thanks $\ldots$''. Instead, try ``R. B. G. thanks$\ldots$''. Put sponsor 
% acknowledgments in the unnumbered footnote on the first page.

% \section*{References}

% Please number citations consecutively within brackets \cite{b1}. The 
% sentence punctuation follows the bracket \cite{b2}. Refer simply to the reference 
% number, as in \cite{b3}---do not use ``Ref. \cite{b3}'' or ``reference \cite{b3}'' except at 
% the beginning of a sentence: ``Reference \cite{b3} was the first $\ldots$''

% Number footnotes separately in superscripts. Place the actual footnote at 
% the bottom of the column in which it was cited. Do not put footnotes in the 
% abstract or reference list. Use letters for table footnotes.

% Unless there are six authors or more give all authors' names; do not use 
% ``et al.''. Papers that have not been published, even if they have been 
% submitted for publication, should be cited as ``unpublished'' \cite{b4}. Papers 
% that have been accepted for publication should be cited as ``in press'' \cite{b5}. 
% Capitalize only the first word in a paper title, except for proper nouns and 
% element symbols.

% For papers published in translation journals, please give the English 
% citation first, followed by the original foreign-language citation \cite{b6}.

% \begin{thebibliography}{00}
% \bibitem{b1} G. Eason, B. Noble, and I. N. Sneddon, ``On certain integrals of Lipschitz-Hankel type involving products of Bessel functions,'' Phil. Trans. Roy. Soc. London, vol. A247, pp. 529--551, April 1955.
% \bibitem{b2} J. Clerk Maxwell, A Treatise on Electricity and Magnetism, 3rd ed., vol. 2. Oxford: Clarendon, 1892, pp.68--73.
% \bibitem{b3} I. S. Jacobs and C. P. Bean, ``Fine particles, thin films and exchange anisotropy,'' in Magnetism, vol. III, G. T. Rado and H. Suhl, Eds. New York: Academic, 1963, pp. 271--350.
% \bibitem{b4} K. Elissa, ``Title of paper if known,'' unpublished.
% \bibitem{b5} R. Nicole, ``Title of paper with only first word capitalized,'' J. Name Stand. Abbrev., in press.
% \bibitem{b6} Y. Yorozu, M. Hirano, K. Oka, and Y. Tagawa, ``Electron spectroscopy studies on magneto-optical media and plastic substrate interface,'' IEEE Transl. J. Magn. Japan, vol. 2, pp. 740--741, August 1987 [Digests 9th Annual Conf. Magnetics Japan, p. 301, 1982].
% \bibitem{b7} M. Young, The Technical Writer's Handbook. Mill Valley, CA: University Science, 1989.
% \end{thebibliography}
% \vspace{12pt}
% \color{red}
% IEEE conference templates contain guidance text for composing and formatting conference papers. Please ensure that all template text is removed from your conference paper prior to submission to the conference. Failure to remove the template text from your paper may result in your paper not being published.

\end{document}
